\documentclass[11pt]{article}

\usepackage[margin=1in]{geometry}
\usepackage{amsmath,amssymb}
\usepackage{array,longtable}
\usepackage[hidelinks]{hyperref}

\setlength{\parindent}{0pt}
\setlength{\parskip}{0.6em}
\renewcommand{\arraystretch}{1.18}

\title{7/24 meeting}
\author{}
\date{}

\begin{document}
\maketitle

\section*{Summary}

\begin{enumerate}
\item \textbf{Missing cells.}
  \(P_I\otimes P_J\) is identified from the observed worker and firm
  marginals, but
  \(\operatorname{Var}_{P_I\otimes P_J}(m_{IJ})\) is generally not:
  it evaluates \(m_{ij}\) at product-supported worker--firm pairs that need
  never be observed.  Hence \(C^w_{\mathrm{assign}}(m)\) is not point
  identified without a valid completion restriction.  BS20 does not recover
  this object; it identifies an observed-match wage-type covariance and
  correlation.
\item \textbf{BS20 versus BS24.}
  The BS20 estimand is a population correlation of worker and firm
  expected-wage types; its estimator is a repeated-match cross-product
  statistic.  Its consistency insight is that distinct-match products remove
  own-observation noise before aggregation.  BS24 does not replace this
  estimator.  It supplies a structural search, acceptance, and bargaining
  mechanism that distinguishes accepted-match effects from meeting effects.
\item \textbf{BS and the project.}
  Both use the same observed matched-wage interface
  \((m,P^{\mathrm{obs}}_{IJ})\), but apply different functionals to it.  The
  later BS Pareto result is exactly additive for \emph{level wages}; the
  project's log-wage surface requires a separate finite-rank or approximation
  bridge.
\item \textbf{Other models.}
  The BS population measure extends whenever a model produces an observed
  matched-wage distribution with finite nondegenerate type moments.  The BS20
  estimator and the later BS meeting/acceptance interpretation require
  additional, model-specific conditions.  Kline, BLM, and GKLP are worked out
  below; AKM/KSS, Tukey/Crippa, Shimer--Smith, and Eeckhout--Kircher follow as
  stated corollaries.
\item \textbf{Exogenous mobility and additive separability.}
  Pure additive separability is an outcome restriction and mobility
  exogeneity is a selection restriction; neither implies the other.
  Sorkin--Warwar obtain nesting because their \emph{additive separability}
  assumption also includes transitory-error exogeneity.  Kline separates the
  statistical AKM condition, causal edge identification, and no selection on
  treatment gains.  Ordinary OLS has the same outcome-versus-selection
  distinction.
\end{enumerate}

Each executive statement is proved below; every sourced claim is followed by
its paper/page reference, and every new equality not taken from a source is
marked as derived.

\section{Question 1: Is \(C^w_{\mathrm{assign}}\) identified, and does BS solve
the missing-cell problem?}


The tables below are the single notation reference for the note.  Moving them
here shortens the advisor-facing argument without changing any symbol or
definition.  A row defining a family such as \(m_{ij}\) also covers indexed
instances such as \(m_{i_0j_0}\), \(m_{ik}\), and
\(m_{g(i),h(j)}\); those are evaluations of the same object, not new
variables.  Primes on indices (\(i',j',m',n',x',y'\)) denote independent
dummy indices in a sum or leave-out construction.  Standard arithmetic,
probability, and linear-algebra operators are not variables.

\subsection{Task 1 notation}

\begin{longtable}{p{0.19\textwidth}p{0.53\textwidth}p{0.21\textwidth}}
\textbf{Symbol} & \textbf{Meaning} & \textbf{Source or status} \\
\hline
\(\mathcal I,\mathcal J\) & Sets of workers and firms. & LOO-v2, p.\ 9. \\
\(I,J\) & Random worker and firm identities inside probabilities,
expectations, variances, and covariances. &
LOO-v2, pp.\ 9--10. \\
\(i,j\) & Particular worker and firm indices.  The pair
\((i_0,j_0)\) denotes a selected missing cell. & LOO-v2, p.\ 9; the subscript
zero is derived notation. \\
\(\mathcal E\) & Observed bipartite edge set:
\((i,j)\in\mathcal E\) when the pair occurs in the data. & LOO-v2, p.\ 9. \\
\(\mathcal T\) & Declared target set on which the maintained schedule is
defined; often \(\operatorname{supp}(P_I\otimes P_J)\). & LOO-v2, pp.\ 9--10. \\
\(Y\) & Observed log wage. & LOO-v2, p.\ 9; BS20, pp.\ 5, 18--20. \\
\(m_{ij}\) & Conditional mean wage on an observed edge and maintained
counterfactual schedule value on an unobserved target cell. & LOO-v2,
pp.\ 9--10, equations (14)--(15); off-support clarification is derived. \\
\(m_{\mathcal E}\) & Vector of maintained wage cells on observed edges. &
LOO-v2, p.\ 19, equation (53). \\
\(P^{\mathrm{obs}}_{IJ}\) & Observed joint worker--firm assignment
distribution. & LOO-v2, pp.\ 9, 15. \\
\(w_{ij}\) & Cell probability
\(P^{\mathrm{obs}}_{IJ}(I=i,J=j)\).  No matrix symbol is assigned here, so
\(W\) remains available for the BLM firm-class membership matrix in Task 4. &
LOO-v2, p.\ 40; matrix-symbol convention is derived. \\
\(P_I,P_J\) & Worker and firm marginals of \(P^{\mathrm{obs}}_{IJ}\). &
LOO-v2, pp.\ 10, 15. \\
\(p_i,q_j\) & Marginal probabilities
\(p_i=\sum_jw_{ij}\) and \(q_j=\sum_iw_{ij}\). & Derived from the definition
of marginals; LOO-v2, pp.\ 10, 40. \\
\(p,q\) & Vectors collecting \(p_i\) and \(q_j\). & Derived notation. \\
\(P^\times\) & Independent-assignment benchmark
\(P^\times=P_I\otimes P_J\), whose cell probability is \(p_iq_j\). &
LOO-v2, p.\ 10, equation (18). \\
\(P,R,r_{ij}\) & A generic probability distribution \(P\), or its probability
matrix \(R=(r_{ij})\), used when stating formulas that apply to either the
observed or product assignment. & Derived notation. \\
\(\operatorname{supp}(P)\) & Cells receiving positive probability under a
distribution \(P\). & Standard notation used in LOO-v2, p.\ 20. \\
\(\mathbb E_P,\operatorname{Var}_P,\operatorname{Cov}_P\) & Expectation,
variance, and covariance under distribution \(P\). & Standard operators. \\
\(\bar m_\times\) & Product-population mean
\(\mathbb E_{P^\times}[m_{IJ}]\). & Derived notation. \\
\(C^w_{\mathrm{assign}}(m)\) & One half of the difference between the observed
and product-assignment variances of the same fixed wage schedule. &
LOO-v2, p.\ 15, equation (42). \\
\(R_{\mathcal E}\) & Selection matrix that extracts observed cells from the
vectorized full schedule. & LOO-v2, p.\ 19, equation (53). \\
\(\mathcal K_{\mathcal E}\) & Missing-cell perturbation space
\(\ker(R_{\mathcal E})\). & LOO-v2, p.\ 19, equation (54). \\
\(e_{i_0j_0}\) & Coordinate vector equal to one at cell \((i_0,j_0)\) and zero
elsewhere. & Derived notation; LOO-v2 uses the same construction on p.\ 39. \\
\(\delta\) & Arbitrary scalar change to one missing cell. & Derived notation. \\
\(m^\delta\) & Observationally equivalent schedule
\(m+\delta e_{i_0j_0}\). & Derived notation. \\
\(a_i,b_j,h_{ij}\) & Product-weighted worker component, firm component, and
double-centered interaction in
\(m_{ij}=\mu+a_i+b_j+h_{ij}\). & LOO-v2, p.\ 11, equations (19)--(23). \\
\(\mu\) & Grand product-weighted mean in the canonical decomposition. &
LOO-v2, p.\ 11, equation (19). \\
\(\mathcal X,\mathcal Y\) & BS20 worker- and firm-characteristic spaces. &
BS20, pp.\ 4--6. \\
\(x,y,z\) & BS20 worker characteristics, firm characteristics, and
match-specific wage determinant; \(x_i\) and \(y_j\) are the characteristics
of worker \(i\) and firm \(j\). & BS20, pp.\ 4--6. \\
\(w(x,y,z)\) & BS20 log-wage function. & BS20, pp.\ 4--6. \\
\(\Phi_x,\Psi_y\) & Conditional distributions of the firm matched to worker
characteristic \(x\), and of the worker matched to firm characteristic \(y\).
& BS20, pp.\ 4--6. \\
\(\lambda(x),\mu(f)\) & BS20 worker and firm types: expected log wage received
conditional on worker characteristic \(x\), and expected log wage paid
conditional on firm characteristic \(f\). &
BS20, p.\ 5, equations (1)--(2). \\
\(\lambda_i^{\mathrm{BS}},\mu_j^{\mathrm{BS}}\) & BS worker and firm
expected-wage types for individual worker \(i\) and firm \(j\).  The
superscript prevents confusion with the project grand mean \(\mu\). &
Derived individual shorthand for BS20, p.\ 5. \\
\(\bar w\) & Common cross-sectional mean log wage. & BS20, p.\ 6,
equation (4). \\
\(\sigma_\lambda,\sigma_\mu\) & Cross-sectional standard deviations of BS20
worker and firm types; their squares are the corresponding variances. &
BS20, p.\ 6, equation (5); p.\ 20, equation (16). \\
\(c_{\mathrm{BS}}\) & Covariance between
\(\lambda(\mathsf X)\) and \(\mu(\mathsf F)\) among
observed matches. & BS20, p.\ 6, equation (6); pp.\ 20--21, equation (17). \\
\(\rho_{\mathrm{BS}}\) & BS20 correlation
\(c_{\mathrm{BS}}/(\sigma_\lambda\sigma_\mu)\). & BS20, p.\ 5,
equation (3). \\
\(M_i,N_j\) & Numbers of observed matches for worker \(i\) and firm \(j\). &
BS20, pp.\ 19, 21. \\
\(m,m',n,n'\) & BS20 focal and other-match indices from the worker and firm
perspectives.  These are indices, not the wage schedule \(m_{ij}\). &
BS20, pp.\ 19--20. \\
\(n(i,m)\) & Firm-side index of the same match that worker \(i\) labels as
match \(m\). & BS20, pp.\ 19--20. \\
\(w^w_{i,m},w^f_{j,n}\) & Average log wage in a match, indexed from the worker
and firm perspectives. & BS20, p.\ 19. \\
\(t^w_{i,m},t^f_{j,n}\) & Durations of those matches. & BS20, p.\ 19. \\
\(T_i^w,T_j^f\) & Total observed employment time of worker \(i\), and total
observed employment supplied by firm \(j\). & BS20, p.\ 19. \\
\(\varepsilon^w_{i,m},\varepsilon^f_{j,n}\) & Worker-side and firm-side
match-wage errors in the BS20 auxiliary equations. & BS20, p.\ 21,
Auxiliary Assumptions 1--4. \\
\(\widehat\lambda_i^{\mathrm{BS}},
\widehat{(\lambda_i^{\mathrm{BS}})^2}\) & BS20-style estimators of a
worker's type and squared type. & BS20, p.\ 22, equation (18); superscript
is the note's canonical clarification. \\
\(\widehat\mu_j^{\mathrm{BS}},
\widehat{(\mu_j^{\mathrm{BS}})^2}\) & BS20-style estimators of a firm's type
and squared type. & BS20, p.\ 23, equation (21); superscript is the note's
canonical clarification. \\
\(\widehat c_{i,m}\) & Leave-focal-match-out product estimating the type
product for the worker and firm in focal match \((i,m)\). & BS20, p.\ 23,
equation (23). \\
\(\widehat{\bar w},\widehat c_{\mathrm{BS}},\widehat\rho_{\mathrm{BS}}\) &
BS20 estimators of the mean wage, type covariance, and type correlation. &
BS20, pp.\ 23--24, equations (19), (24), and (25). \\
\(\tau\) & BS20 replication index: the numbers of each worker and firm
characteristic are multiplied by \(\tau\), and consistency is studied as
\(\tau\to\infty\). & BS20, p.\ 18. \\
\(s_{ij}\) & Constructed additive BS type-score schedule
\(\lambda_i^{\mathrm{BS}}+\mu_j^{\mathrm{BS}}-\bar w\). & Derived in
(T1.10). \\
\end{longtable}
\subsection{1.1 Objects and answer}

Let \(\mathcal E\subseteq\mathcal I\times\mathcal J\) be the observed
worker--firm edge set.  On an observed edge,
\[
m_{ij}=\mathbb E[Y\mid I=i,J=j],
\qquad
w_{ij}=P^{\mathrm{obs}}_{IJ}(i,j).
\tag{T1.1}
\]
On a target cell outside \(\mathcal E\), \(m_{ij}\) denotes the maintained
counterfactual value of the same wage schedule; it is not an observed
conditional mean.  This distinction is the source of the identification
problem below.
The observed marginals and product reference law are
\[
p_i=\sum_jw_{ij},
\qquad
q_j=\sum_iw_{ij},
\qquad
P^\times(i,j)=p_iq_j.
\tag{T1.2}
\]
Thus the product \emph{measure} is observable.  The problem is its integrand:
\[
C^w_{\mathrm{assign}}(m)
=
\frac12\left\{
\operatorname{Var}_{P^{\mathrm{obs}}_{IJ}}(m_{IJ})
-\operatorname{Var}_{P^\times}(m_{IJ})
\right\}.
\tag{T1.3}
\]
The first variance uses only observed edges because \(w_{ij}=0\) off
\(\mathcal E\).  The second uses every pair with \(p_iq_j>0\), including
missing cells.

\textit{Source: LOO-v2, pp.\ 9--11 and 15--20.}

\subsection{1.2 Missing-cell nonidentification proof}

Assume that some \((i_0,j_0)\notin\mathcal E\) satisfies
\(p_{i_0}q_{j_0}>0\).  Construct two schedules that agree on every observed
cell:
\[
m^\delta_{ij}
=
\begin{cases}
m_{ij}+\delta,&(i,j)=(i_0,j_0),\\
m_{ij},&\text{otherwise}.
\end{cases}
\tag{T1.4}
\]
Because the changed cell has zero observed probability,
\[
\operatorname{Var}_{P^{\mathrm{obs}}_{IJ}}(m^\delta_{IJ})
=
\operatorname{Var}_{P^{\mathrm{obs}}_{IJ}}(m_{IJ}).
\tag{T1.5}
\]
Let \(\bar m_\times=\mathbb E_{P^\times}[m_{IJ}]\).  Because
\(m^\delta\) changes only one cell, its product-law mean and second moment are
\[
\mathbb E_{P^\times}[m^\delta_{IJ}]
=
\bar m_\times+p_{i_0}q_{j_0}\delta,
\tag{T1.5a}
\]
\[
\mathbb E_{P^\times}[(m^\delta_{IJ})^2]
=
\mathbb E_{P^\times}[m_{IJ}^2]
+p_{i_0}q_{j_0}
\left(2\delta m_{i_0j_0}+\delta^2\right).
\tag{T1.5b}
\]
Subtracting the old variance from the new variance, using
\(\operatorname{Var}(X)=\mathbb E[X^2]-\mathbb E[X]^2\), gives
\[
\begin{aligned}
&\operatorname{Var}_{P^\times}(m^\delta_{IJ})
-\operatorname{Var}_{P^\times}(m_{IJ})\\
&\qquad =
p_{i_0}q_{j_0}(2\delta m_{i_0j_0}+\delta^2)
-\left[
(\bar m_\times+p_{i_0}q_{j_0}\delta)^2-\bar m_\times^2
\right]\\
&\qquad =
2p_{i_0}q_{j_0}\delta(m_{i_0j_0}-\bar m_\times)
+p_{i_0}q_{j_0}(1-p_{i_0}q_{j_0})\delta^2.
\end{aligned}
\tag{T1.6}
\]
Because \(0<p_{i_0}q_{j_0}<1\), the right side of (T1.6) is a nonconstant
quadratic in \(\delta\).  Hence there exists \(\delta\neq0\) for which it is
nonzero.  Combining (T1.3), (T1.5), and (T1.6), the two schedules then
generate identical observed data but different
\(C^w_{\mathrm{assign}}\).  This proves nonidentification on the unrestricted
incomplete schedule.  The factor \(1/2\) in (T1.3) changes the scale of the
difference, not this conclusion.

Repeated observations of the same observed edges improve estimation of
\(m_{ij}\) on \(\mathcal E\), but they do not reveal
\(m_{i_0j_0}\).  The BS minimum-match requirement therefore does not alter
this proof.

\textit{Source: LOO-v2, pp.\ 19--20, Proposition 4 and Corollary 1; the
single-cell expansion is derived.}

\subsection{1.3 What BS20 identifies instead}

BS20 defines observed-match expected-wage types
\[
\lambda_i^{\mathrm{BS}}
=
\sum_jP^{\mathrm{obs}}(J=j\mid I=i)m_{ij},
\qquad
\mu_j^{\mathrm{BS}}
=
\sum_iP^{\mathrm{obs}}(I=i\mid J=j)m_{ij}.
\tag{T1.7}
\]
Its covariance and correlation are
\[
c_{\mathrm{BS}}
=
\operatorname{Cov}_{P^{\mathrm{obs}}_{IJ}}
(\lambda_I^{\mathrm{BS}},\mu_J^{\mathrm{BS}}),
\qquad
\rho_{\mathrm{BS}}
=
\frac{c_{\mathrm{BS}}}{\sigma_\lambda\sigma_\mu}.
\tag{T1.8}
\]
These objects require conditional means under the observed assignment law,
not completion of every product-supported cell.

They are not the project product-reference main effects.  For example,
\[
a_i=\sum_jq_jm_{ij}-\mu,
\qquad
\mu=\sum_{i,j}p_iq_jm_{ij}
\tag{T1.9}
\]
uses the unconditional firm marginal \(q_j\), whereas
\(\lambda_i^{\mathrm{BS}}\) uses
worker \(i\)'s observed conditional firm distribution.  Under dependent
assignment the weights differ, and (T1.9) can require missing cells.

The exact bridge is obtained by replacing the original wage surface with the
additive BS score
\[
s_{ij}=\lambda_i^{\mathrm{BS}}+\mu_j^{\mathrm{BS}}-\bar w.
\tag{T1.10}
\]
The observed and product laws have the same worker and firm marginals.
Therefore the marginal variance terms cancel, while the product-law
covariance is zero:
\[
\begin{aligned}
\operatorname{Var}_{P^{\mathrm{obs}}}(s_{IJ})
&=
\sigma_\lambda^2+\sigma_\mu^2+2c_{\mathrm{BS}},\\
\operatorname{Var}_{P^\times}(s_{IJ})
&=
\sigma_\lambda^2+\sigma_\mu^2.
\end{aligned}
\tag{T1.10a}
\]
Substitution into the definition of the assignment contribution gives
\[
\begin{aligned}
C^w_{\mathrm{assign}}(s)
&=
\frac12\left\{
\operatorname{Var}_{P^{\mathrm{obs}}}(s_{IJ})
-\operatorname{Var}_{P^\times}(s_{IJ})
\right\}\\
&=
\frac12\left\{
2\operatorname{Cov}_{P^{\mathrm{obs}}}
(\lambda_I^{\mathrm{BS}},\mu_J^{\mathrm{BS}})
\right\}\\
&=c_{\mathrm{BS}}.
\end{aligned}
\tag{T1.11}
\]
Hence BS gives a valid observed-support sorting statistic, or the assignment
contribution of the \emph{replacement} schedule \(s\).  It equals the
original \(C^w_{\mathrm{assign}}(m)\) only under an additional bridge
restriction, such as a correctly specified completion model tying \(m_{ij}\)
to identified BS types.  Substituting \(c_{\mathrm{BS}}\) for the original
target without such a restriction changes the estimand.

\textit{Source: BS20, pp.\ 5--7 and 18--24.  Equations (T1.9)--(T1.11) and
the non-equivalence statement are derived.}

\subsection{1.4 Operational conclusion}

\begin{longtable}{p{0.30\textwidth}p{0.17\textwidth}p{0.43\textwidth}}
\textbf{Object} & \textbf{Identified?} & \textbf{Reason} \\
\hline
\(P_I\otimes P_J\) & Yes & Its factors are observed assignment marginals. \\
\(m_{ij}\) on \(\mathcal E\) & Estimable & Requires adequate wage
observations on the observed edge. \\
\(\operatorname{Var}_{P^\times}(m_{IJ})\) & Generally no & Requires
\(m_{ij}\) in product-supported missing cells. \\
\(C^w_{\mathrm{assign}}(m)\) & Generally no & Contains the unidentified
product variance. \\
\((c_{\mathrm{BS}},\rho_{\mathrm{BS}})\) & Yes under BS sampling assumptions &
Uses observed-match wage types and repeated-match cross-products. \\
\end{longtable}

\section{Question 2: What are the BS estimand and estimator, why is the
estimator consistent, and what changes in BS24?}

\subsection{Task 2 notation}

The chart retains all Task 1 notation.  The meeting rate is written
\(\rho_{\mathrm{meet}}\) here to prevent confusion with the BS20 sorting
correlation \(\rho_{\mathrm{BS}}\).  BS24 writes the meeting-rate parameter as
\(\rho\).

\begin{longtable}{p{0.20\textwidth}p{0.52\textwidth}p{0.21\textwidth}}
\textbf{Symbol} & \textbf{Meaning} & \textbf{Source or status} \\
\hline
\(\mathsf X,\mathsf F,Y;x,f;\mathsf X_i,\mathsf F_j\) & BS20 random worker
characteristic, random firm
characteristic, observed log wage, realizations of the two characteristics,
and the individual worker/firm realizations.  The letter \(F\) avoids reusing
the wage symbol \(Y\) for a firm characteristic. & BS20, pp.\ 4--6;
relabeling is derived. \\
\(x,y\) & Ordered worker and firm types in BS24's structural model. & BS24,
pp.\ 4--5. \\
\(\lambda(x),\mu(f)\) & BS20 expected log-wage worker and firm types.  These
remain the canonical symbols for the BS20 estimand. & BS20, p.\ 5,
equations (1)--(2). \\
\(\bar w\) & BS20 common mean log wage.  This scalar is distinct from the
BS24 type-specific reservation wage \(\bar w_x\). & BS20, p.\ 6,
equation (4). \\
\(\sigma_\lambda^2,\sigma_\mu^2\) & Variances of BS20 worker and firm wage
types. & BS20, pp.\ 6, 20, equations (5) and (16). \\
\(c_{\mathrm{BS}}\) & Covariance of BS20 wage types in observed matches. &
BS20, pp.\ 6, 20--21, equations (6) and (17). \\
\(\rho_{\mathrm{BS}}\) & BS20 sorting estimand
\(c_{\mathrm{BS}}/(\sigma_\lambda\sigma_\mu)\). & BS20, p.\ 5,
equation (3). \\
\(\rho_{\mathrm{AKM}}\) & Correlation between AKM worker and firm effects in
observed matches, used by BS20 as a comparator. & BS20, pp.\ 7--9. \\
\(\widehat{\bar w}\) & BS20 estimator of the common mean wage. & BS20,
p.\ 23, equation (19). \\
\(\widehat\sigma_\lambda^2,\widehat\sigma_\mu^2\) & BS20 estimators of the
two type variances. & BS20, pp.\ 23--24, equations (20) and (22). \\
\(\widehat c_{\mathrm{BS}},\widehat\rho_{\mathrm{BS}}\) & BS20 covariance and
correlation estimators. & BS20, p.\ 24, equations (24)--(25). \\
\(M_i,N_j\) & Numbers of distinct observed matches for worker \(i\) and firm
\(j\). & BS20, pp.\ 19--22. \\
\(m,m',n,n'\) & Focal and other-match indices in BS20.  They are not the
project wage schedule \(m_{ij}\). & BS20, pp.\ 19--23. \\
\(j(i,m),n(i,m)\) & Firm in worker \(i\)'s focal match \(m\), and the index
of that same match from the firm's side. & Derived clarification of BS20,
p.\ 23, equation (23). \\
\(w^w_{i,m},w^f_{j,n}\) & Match-average log wages indexed from the worker
and firm sides. & BS20, p.\ 19. \\
\(\varepsilon^w_{i,m},\varepsilon^f_{j,n}\) & Mean-zero auxiliary
match-wage errors. & BS20, pp.\ 21--23, Auxiliary Assumptions 1--4. \\
\(\widehat c_{i,m}^{\,\mathrm{raw}}\) & Product of worker-side and firm-side
leave-focal-match-out average wages in (T2.5b). & Derived notation for BS20,
p.\ 23, equation (23). \\
\(\tau\) & BS20 replication index governing the number of workers and firms;
the number of observations per unit may remain bounded as
\(\tau\to\infty\). & BS20, p.\ 18. \\
\(f_{x,y}\) & BS24 deterministic production scale for worker type \(x\) and
firm type \(y\). & BS24, pp.\ 5--7. \\
\(z,G,g\) & Match-specific productivity, its cumulative distribution, and
its density. & BS24, pp.\ 5--6. \\
\(z_0,\theta\) & Pareto scale and shape parameters.  The exact finite-mean
results use \(\theta>1\). & BS24, pp.\ 10--12. \\
\(\gamma\) & Worker's Nash-bargaining share in the accepted-wage equation. &
BS24, pp.\ 4--8. \\
\(\rho_{\mathrm{meet}}\) & Canonical symbol for the BS24 meeting-rate
parameter, written \(\rho\) in that paper. & BS24, p.\ 5. \\
\(u_x^{\mathrm{BS}},v_y^{\mathrm{BS}}\) & Steady-state measures of
unemployed type-\(x\) workers and vacant type-\(y\) jobs.  The paper omits
the superscripts. & BS24, pp.\ 5, 8--9; superscripts are the note's
canonical clarification. \\
\(\bar w_x,\bar\pi_y\) & Reservation wage of a type-\(x\) worker and
reservation profit of a type-\(y\) firm. & BS24, pp.\ 6--8. \\
\(\bar z_{x,y}\) & Match-acceptance productivity threshold
\((\bar w_x+\bar\pi_y)/f_{x,y}\). & BS24, p.\ 7, equation (6). \\
\(p_{x,y}\) & Canonical shorthand for the match-acceptance probability
\(1-G(\bar z_{x,y})\). & Derived from BS24 equation (6). \\
\(\mathcal G_{x,y}\) & Worker value gain created by a potential
\((x,y)\) meeting, defined to be zero when the meeting is rejected. & Derived
notation for BS24/25 revision, pp.\ 9--11. \\
\(V^s_{x,y}(z)\) & Total surplus created by an \((x,y,z)\) match. & BS24,
pp.\ 6--7, equation (5). \\
\(W_{x,y}(z)\) & Equilibrium wage paid in an accepted
\((x,y,z)\) match. & BS24, pp.\ 7--8, equation (9). \\
\(\bar\phi_{x,y}\) & Steady-state measure of accepted matches between types
\(x\) and \(y\). & BS24, p.\ 8, equation (13). \\
\(\omega_{i,j}\) & Observed wage of worker \(i\) at firm \(j\) in the BS24/25
analytic wage equation. & BS24/25 revision, pp.\ 10--12. \\
\(\alpha_x^*,\psi_y^*\) & BS24/25 type-level worker and firm components of
accepted-match wages. & BS24/25 revision, pp.\ 11--13,
equations (21)--(22). \\
\(\varepsilon_{i,j}\) & Conditional mean-zero accepted-match wage residual in
the Pareto model. & BS24/25 revision, pp.\ 11--13,
equations (21)--(23). \\
\(\bar V^s_{x,y}\) & BS24 average surplus conditional on an \((x,y)\) meeting
becoming a match. & BS24, pp.\ 31--32, equation (33). \\
\(\bar V^{s,m}_{x,y}\) & BS24 expected surplus from a random \((x,y)\)
meeting before match quality is known.  The superscript \(m\) is the paper's
notation for ``meeting,'' not the project wage schedule. & BS24, pp.\ 31--32,
equation (32). \\
\(\operatorname{ATT}_x(y)\) & Worker-value gain from an \((x,y)\) meeting
conditional on acceptance in the 2025 revision. & BS24/25 revision,
pp.\ 9--11, equation (19). \\
\(\operatorname{ITT}_x(y)\) & Worker-value gain from creating a random
\((x,y)\) meeting before knowing whether it will be accepted. & BS24/25
revision, pp.\ 9--11, equations (17)--(18). \\
\(\alpha_i,\psi_j\) & Individual worker and firm fixed effects in the
conventional two-way fixed-effects regression.  These are estimators'
parameters, distinct from the type-level structural components
\(\alpha_x^*,\psi_y^*\). & BS24/25 revision, pp.\ 15--17,
equation (30). \\
\(\widehat\alpha_i,\widehat\psi_j\) & OLS estimates of the individual
two-way fixed effects. & BS24/25 revision, pp.\ 16--18. \\
\(B\) & Multiplicative attenuation factor for estimated firm-component
differences under endogenous separations in the 2025 revision. & BS24/25
revision, pp.\ 26--28, Proposition 7. \\
\end{longtable}
\textbf{Version convention.}
BS20 is \emph{High Wage Workers Work for High Wage Firms}.  BS24 denotes the
November 2024 Richmond Fed/NBER working paper; despite the shorthand used in
the project discussion, it should not be described as a published QJE
article.  ``BS24/25 revision'' denotes the later local revision used for the
ITT/ATT terminology and added results.
\textit{Sources: paper title pages and the authors' version information.}

\subsection{2.1 BS20: object, estimand, estimator, and estimate}

Let \(\mathsf X\) denote a worker characteristic, \(\mathsf F\) a firm
characteristic, and
\(Y\) the observed log wage.  This keeps \(Y\) as the wage outcome throughout
the note.  BS20 defines
\[
\lambda(x)=\mathbb E[Y\mid \mathsf X=x],
\qquad
\mu(f)=\mathbb E[Y\mid \mathsf F=f].
\tag{T2.1}
\]
For an individual worker and firm,
\(\lambda_i^{\mathrm{BS}}=\lambda(\mathsf X_i)\) and
\(\mu_j^{\mathrm{BS}}=\mu(\mathsf F_j)\).
The population object of interest and estimand is
\[
\rho_{\mathrm{BS}}
=
\operatorname{Corr}(\lambda(\mathsf X),\mu(\mathsf F)).
\tag{T2.2}
\]
The covariance \(c_{\mathrm{BS}}\), the two type variances, and the mean wage
are component estimands.  The estimator
\(\widehat\rho_{\mathrm{BS}}\) is the random sample statistic constructed
from estimators of those four moments.  Its realized numerical value is the
estimate.  The AKM correlation is a comparator, not the definition of
\(\rho_{\mathrm{BS}}\).

\begin{longtable}{p{0.22\textwidth}p{0.27\textwidth}p{0.42\textwidth}}
\textbf{Layer} & \textbf{Object} & \textbf{Meaning} \\
\hline
Economic concept & Vertical wage sorting & High expected-wage workers match
with high expected-wage firms. \\
Estimand & \(\rho_{\mathrm{BS}}\) & Fixed population functional. \\
Component estimands & \(\bar w,\sigma_\lambda^2,\sigma_\mu^2,c_{\mathrm{BS}}\)
& Moments entering the correlation. \\
Estimator & \(\widehat\rho_{\mathrm{BS}}\) & Random statistic formed from
observed wages and match histories. \\
Estimate & Realized \(\widehat\rho_{\mathrm{BS}}\) & Dataset-specific number. \\
\end{longtable}

\textit{Source: BS20, pp.\ 5--7, equations (1)--(6).}

\subsection{2.2 The consistency insight}

The hard case is a short panel: the number \(M_i\) of distinct matches for a
worker may remain small, so a worker-specific sample mean need not consistently
estimate \(\lambda_i^{\mathrm{BS}}\).  Write a wage from distinct worker
matches as
\[
w^w_{i,m}=\lambda_i^{\mathrm{BS}}+\varepsilon^w_{i,m},
\qquad
\mathbb E[\varepsilon^w_{i,m}\mid\lambda_i^{\mathrm{BS}}]=0.
\tag{T2.3}
\]
Assume \(M_i\geq2\).  For \(m\neq m'\), conditional independence of
distinct-match errors yields
\[
\begin{aligned}
\mathbb E[w^w_{i,m}w^w_{i,m'}\mid\lambda_i^{\mathrm{BS}}]
&=
(\lambda_i^{\mathrm{BS}})^2
+\lambda_i^{\mathrm{BS}}
\mathbb E[\varepsilon^w_{i,m}+\varepsilon^w_{i,m'}
\mid\lambda_i^{\mathrm{BS}}]\\
&\quad+
\mathbb E[\varepsilon^w_{i,m}\varepsilon^w_{i,m'}
\mid\lambda_i^{\mathrm{BS}}]\\
&=(\lambda_i^{\mathrm{BS}})^2.
\end{aligned}
\tag{T2.4}
\]
The two linear terms are zero by conditional mean zero; conditional
independence makes the error product's conditional expectation equal the
product of those two zero means.
Therefore the U-statistic
\[
\widehat{(\lambda_i^{\mathrm{BS}})^2}
=
\frac{1}{M_i(M_i-1)}
\sum_{m\neq m'}w^w_{i,m}w^w_{i,m'}
\tag{T2.5}
\]
is unbiased for \((\lambda_i^{\mathrm{BS}})^2\) even if \(M_i\) is fixed.
For a firm with \(N_j\geq2\), the corresponding estimator is
\[
\widehat{(\mu_j^{\mathrm{BS}})^2}
=
\frac{1}{N_j(N_j-1)}
\sum_{n\neq n'}w^f_{j,n}w^f_{j,n'}.
\tag{T2.5a}
\]
Expanding the product as in (T2.4), now under the firm-side auxiliary
equation, proves that (T2.5a) is unbiased for
\((\mu_j^{\mathrm{BS}})^2\).

For \(c_{\mathrm{BS}}\), let \(j(i,m)\) be worker \(i\)'s firm in focal match
\(m\), and let \(n(i,m)\) be that same match's index from firm \(j(i,m)\)'s
side.  Define
\[
\widehat c_{i,m}^{\,\mathrm{raw}}
=
\left[
\frac{1}{M_i-1}\sum_{m'\neq m}w^w_{i,m'}
\right]
\left[
\frac{1}{N_{j(i,m)}-1}
\sum_{n'\neq n(i,m)}w^f_{j(i,m),n'}
\right].
\tag{T2.5b}
\]
Auxiliary Assumption 3 makes the two leave-out averages independent
conditional on the focal types.  Each average is separately unbiased, so
\[
\mathbb E[
\widehat c_{i,m}^{\,\mathrm{raw}}
\mid\lambda_i^{\mathrm{BS}},\mu_{j(i,m)}^{\mathrm{BS}}]
=
\lambda_i^{\mathrm{BS}}\mu_{j(i,m)}^{\mathrm{BS}}.
\tag{T2.5c}
\]
Duration-weighted aggregation over focal matches estimates
\(\mathbb E[\lambda_I^{\mathrm{BS}}\mu_J^{\mathrm{BS}}]\); subtracting the
square of the estimated common mean gives \(c_{\mathrm{BS}}\).

The logical order is:
\[
\text{distinct-match products}
\Longrightarrow
\text{unbiased scores for latent moments}
\Longrightarrow
\text{law of large numbers across units/matches}
\Longrightarrow
\widehat\rho_{\mathrm{BS}}\overset p\longrightarrow\rho_{\mathrm{BS}}.
\tag{T2.6}
\]
The final continuous-mapping step also requires
\(\sigma_\lambda^2>0\) and \(\sigma_\mu^2>0\).
The theorem therefore identifies an aggregate functional without requiring
consistent estimation of every individual latent type.  Its assumptions
concern stable types, distinct-match error independence, duration weighting,
cross-side leave-out independence, and cross-unit independence sufficient for
the law of large numbers.
They are not implied merely by observing a connected mobility graph.

\textit{Source: BS20, pp.\ 18--24, Auxiliary Assumptions 1--4 and
Propositions 2--5.}

\subsection{2.3 BS24: a structural continuation, not a new estimator}

BS24 starts before a match exists.  A worker type \(x\) and firm type \(y\)
meet, draw match productivity \(z\), and accept when
\[
z\geq\bar z_{x,y}
=
\frac{\bar w_x+\bar\pi_y}{f_{x,y}},
\qquad
p_{x,y}=1-G(\bar z_{x,y}).
\tag{T2.7}
\]
Conditional on acceptance, Nash bargaining determines the wage.  The model
therefore generates production, meetings, acceptance, wages, and the observed
match distribution.

If \(\mathcal G_{x,y}\) denotes the worker gain from a potential meeting and rejected
meetings yield zero gain, define
\[
\operatorname{ITT}_x(y)
=\mathbb E[\mathcal G_{x,y}\mid x,y,\text{meeting}],
\tag{T2.8}
\]
\[
\operatorname{ATT}_x(y)
=
\mathbb E[\mathcal G_{x,y}\mid x,y,\text{meeting, accepted}].
\tag{T2.9}
\]
The law of total expectation gives
\[
\begin{aligned}
\operatorname{ITT}_x(y)
&=
p_{x,y}\operatorname{ATT}_x(y)
+(1-p_{x,y})
\mathbb E[\mathcal G_{x,y}\mid x,y,\text{meeting, rejected}]\\
&=p_{x,y}\operatorname{ATT}_x(y),
\end{aligned}
\tag{T2.10}
\]
because the rejected-meeting gain is zero by definition.
BS24 explicitly treats the ITT as the policy object of interest: it values a
new meeting before knowing whether a match forms.  The ATT is the selected
accepted-meeting object that the firm wage component can reveal under the
Pareto structure.
Administrative wages reveal accepted matches.  They do not reveal rejected
meetings or \(p_{x,y}\) without structural information.  Thus a wage
fixed-effect contrast may map to an ATT under additional conditions but does
not by itself identify an ITT.

\textit{Source: BS24/25 revision, pp.\ 6--14.}

\subsection{2.4 What BS24 adds, loses, and contributes to this project}

\begin{longtable}{p{0.18\textwidth}p{0.34\textwidth}p{0.39\textwidth}}
\textbf{Dimension} & \textbf{BS20} & \textbf{BS24/25} \\
\hline
Main object & Observed-match wage-type correlation. & Structural production,
selection, bargaining, and match stocks; the policy object is the ITT, while
accepted wages are linked to the ATT. \\
Generality & Several models can generate the same functional. & Sharper
results require ordered types and restrictions such as Pareto shocks,
monotonicity, and log-supermodularity. \\
Empirics & Direct repeated-match estimator and consistency theorem. &
Calibration and structural diagnostics, not a replacement for the BS20
estimator. \\
Main lesson & A portable measure of vertical wage sorting. & Accepted-wage
additivity need not imply additive production; mobility diagnostics need not
remove structural selection. \\
Project use & Core semistructural measurement benchmark. & Structural
interpretation and sensitivity model. \\
\end{longtable}

BS20 should therefore be the project's core BS reference.  BS24 should be
used when interpreting production, selection, or policy.  Imposing BS24's
random-search/Pareto structure on the headline project objects would make the
paper substantially more structural.  Neither paper solves Task 1's original
missing-cell target without an additional bridge.

\textit{Sources: BS20, Sections 3--4; BS24, Sections 2--8; BS24/25 revision,
Sections 2--6.  The comparison is derived.}

\section{Question 3: What is the exact BS--project connection?}

\subsection{Task 3 notation}

This chart inherits all notation from Tasks 1 and 2 without changing it.  It
lists every new symbol used in Task 3 and repeats the old symbols needed to
read the map without looking backward.  In particular,
\(\rho_{\mathrm{meet}}\) remains the meeting rate and
\(\rho_{\mathrm{BS}}\) remains the sorting correlation.

\begin{longtable}{p{0.20\textwidth}p{0.52\textwidth}p{0.21\textwidth}}
\textbf{Symbol} & \textbf{Meaning} & \textbf{Source or status} \\
\hline
\(\mathcal I,\mathcal J;i,j\) & Project worker and firm sets; individual
worker and firm indices. & LOO-v2, p.\ 9. \\
\(I,J\) & Random worker and firm identities when used inside probabilities,
expectations, and variances. & LOO-v2, pp.\ 9--10. \\
\(i',j',k,\ell,x',y'\) & Dummy individual, type, and class indices used
inside sums; a prime distinguishes the dummy from the focal index. & Derived
indexing notation. \\
\(x_i,y_j\) & BS structural types of worker \(i\) and firm \(j\). &
BS24/25 revision, pp.\ 4--5. \\
\(x_1,x_2,y_1,y_2\) & Two ordered worker types and two ordered firm types used
to state the BS log-wage submodularity inequality. & BS24/25 revision,
pp.\ 17--18. \\
\(P^{\mathrm{obs}}_{IJ},w_{ij}\) & Observed individual assignment law and
its cell probabilities. & LOO-v2, pp.\ 9, 40. \\
\(P_I,P_J\) & Worker and firm marginals of the observed individual assignment
law. & LOO-v2, pp.\ 10, 15. \\
\(p_i,q_j\) & Individual worker and firm marginals of
\(P^{\mathrm{obs}}_{IJ}\). & LOO-v2, pp.\ 10, 40. \\
\(P^\times\) & Project product benchmark
\(P_I\otimes P_J\), with cell mass \(p_iq_j\). & LOO-v2, p.\ 10. \\
\(Y,m_{ij}\) & Project log-wage outcome and maintained conditional mean log
wage for pair
\((i,j)\), including target pairs when a completion is imposed. & LOO-v2,
pp.\ 9--10. \\
\(\mu,a_i,b_j,h_{ij}\) & Product-weighted grand mean, worker component, firm
component, and double-centered interaction in
\(m_{ij}=\mu+a_i+b_j+h_{ij}\). & LOO-v2, p.\ 11. \\
\(\mathbf u_i,\mathbf v_j,\Lambda\) & Project worker and firm interaction
factor vectors and coefficient matrix in
\(h_{ij}=\mathbf u_i'\Lambda\mathbf v_j\).  Boldface prevents confusion with
BS unemployment and vacancy measures. & LOO-v2, pp.\ 21--23; boldface is a
notational clarification. \\
\(r_0\) & Maintained project interaction rank in Task 3.  It is distinct from
the BS discount rate \(r\). & Derived clarification of LOO-v2's \(r\). \\
\(Q_F,H_F,A_h,C^w_{\mathrm{assign}}\) & Project average within-worker
cross-firm dispersion, heterogeneity in firm contrasts, mean favorable-cell
assignment, and fixed-schedule variance assignment contribution. & LOO-v2,
pp.\ 12--17. \\
\(\lambda_i^{\mathrm{BS}},\mu_j^{\mathrm{BS}}\) & BS20 expected log-wage type
of individual worker \(i\) and firm \(j\).  Superscripts distinguish the firm
type from the project's scalar \(\mu\). & BS20, p.\ 5; superscripts are
derived clarification. \\
\(r_i^{\mathrm{comp}},s_j^{\mathrm{comp}}\) & Observed counterpart-
composition and interaction terms entering the worker and firm BS types in
(T3.8a). & Derived. \\
\(c_{\mathrm{BS}},\rho_{\mathrm{BS}}\) & BS20 covariance and correlation
between worker and firm wage types in observed matches. & BS20, pp.\ 5--7. \\
\(\bar w,\sigma_\lambda,\sigma_\mu\) & BS20 common observed mean log wage and
standard deviations of its worker and firm wage types.  The scalar
\(\bar w\) is distinct from the later paper's reservation wage
\(\bar w_x\). & BS20, pp.\ 5--7. \\
\(T_{\mathrm{BS}}(m,P^{\mathrm{obs}}_{IJ})\) & Derived name for the mapping
from a wage schedule and assignment law to \(\rho_{\mathrm{BS}}\). & Derived
notation. \\
\(\mathcal S_{\mathrm{interface}}\) & Ordered pair
\((m,P^{\mathrm{obs}}_{IJ})\), used as a name for the common empirical
interface. & Derived notation. \\
\(f_{x,y}\) & BS deterministic type-pair production scale. & BS24/25
revision, pp.\ 5--7. \\
\(z,G,g,g(z),z_0,\theta\) & Match-specific productivity draw, its CDF,
density symbol and density evaluated at \(z\),
and the Pareto scale and shape parameters. & BS24/25 revision, pp.\ 5,
11--12. \\
\(t\) & A possible realized match-quality value used as the upper argument
of the truncated CDF. & Derived dummy value. \\
\(\gamma\) & Worker's Nash-bargaining share in the accepted-wage equation. &
BS24/25 revision, pp.\ 4--8. \\
\(\rho_{\mathrm{meet}}\) & BS meeting-rate parameter, denoted \(\rho\) in
the paper. & BS24/25 revision, p.\ 5; canonical renaming from Task 2. \\
\(u_x^{\mathrm{BS}},v_y^{\mathrm{BS}}\) & Steady-state measures of
unemployed type-\(x\) workers and vacant type-\(y\) jobs.  Superscripts
separate them from project interaction factors. & BS24/25 revision,
pp.\ 5, 8; superscripts are the note's canonical clarification. \\
\(\bar w_x,\bar\pi_y\) & Type-\(x\) reservation wage and type-\(y\)
reservation profit. & BS24/25 revision, pp.\ 6--8. \\
\(\bar z_{x,y},p_{x,y}\) & Acceptance threshold and acceptance probability,
\(\bar z_{x,y}=(\bar w_x+\bar\pi_y)/f_{x,y}\) and
\(p_{x,y}=1-G(\bar z_{x,y})\). & BS24/25 revision, pp.\ 6--8; the \(p\)
shorthand is canonical from Task 2. \\
\(V^s_{x,y}(z),W_{x,y}(z)\) & Total match surplus and equilibrium accepted
wage for a realized \(z\). & BS24/25 revision, pp.\ 6--8. \\
\(\bar\phi_{x,y},\bar\Phi\) & Steady-state mass of accepted \((x,y)\)
matches and total accepted-match mass
\(\bar\Phi=\sum_{x,y}\bar\phi_{x,y}\). & First symbol: BS24/25 revision,
p.\ 8; total is derived. \\
\(\Pi^{\mathrm{obs}}_{x,y}\) & Observed joint type distribution among
accepted matches, \(\bar\phi_{x,y}/\bar\Phi\). & Derived from BS steady-state
match masses. \\
\(\pi_x,\kappa_y\) & Worker- and firm-type marginals of
\(\Pi^{\mathrm{obs}}_{x,y}\).  These are not the reservation profit
\(\bar\pi_y\). & Derived notation. \\
\(\omega_{i,j},\varepsilon_{i,j}\) & Realized accepted wage in levels and
its Pareto-model conditional mean-zero residual. & BS24/25 revision,
pp.\ 11--13. \\
\(\alpha_x^*,\psi_y^*\) & BS Pareto worker- and firm-type components of the
accepted wage in levels. & BS24/25 revision, pp.\ 11--13. \\
\(m^{L}_{ij}\) & Conditional mean accepted wage in levels for individual
pair \((i,j)\). & Derived from BS Proposition 2. \\
\(A_i,B_j,\bar A,\bar B\) & Individual lifts of the BS worker and firm
level-wage components and their product-weighted marginal means. & Derived
notation for the level-wage bridge. \\
\(\mu^L,a_x^L,b_y^L,h_{x,y}^L\) & Canonical project decomposition applied to
the conditional mean accepted wage in levels. & Derived notation. \\
\(Q_F^L,H_F^L,A_h^L,C_{\mathrm{assign}}^{w,L}\) & Project functionals applied
to the level-wage schedule rather than the original log-wage schedule. &
Derived notation. \\
\(\ell_{x,y},m^{\log}_{ij}\) & Conditional mean accepted log wage in a type
cell and its individual-pair lift,
\(m^{\log}_{ij}=\ell_{x_i,y_j}\). & First object: BS24/25 revision,
pp.\ 17--18; individual lift is derived. \\
\(\ell_{x\cdot},\ell_{\cdot y},\ell_{\cdot\cdot}\) & Product-weighted
worker-type row mean, firm-type column mean, and grand mean of
\(\ell_{x,y}\), using \(\kappa_y\) and \(\pi_x\). & Derived canonical
centering. \\
\(M^\ell\) & Double-centered type-cell mean log-wage matrix with element
\(M^\ell_{x,y}=\ell_{x,y}-\ell_{x\cdot}-\ell_{\cdot y}
+\ell_{\cdot\cdot}\). & Derived from LOO-v2 canonical decomposition. \\
\(K_\ell,L_\ell,\Sigma_\ell,R_\ell\) & Rank and thin singular-value
decomposition of \(M^\ell\):
\(K_\ell=\operatorname{rank}(M^\ell)\) and
\(M^\ell=L_\ell\Sigma_\ell R_\ell'\). & Standard linear algebra; derived
application. \\
\(e_x,e_y\) & Coordinate vectors selecting worker-type row \(x\) and
firm-type row \(y\) in the singular-value decomposition. & Standard linear
algebra. \\
\(\mathbf u_x,\mathbf v_y\) & Worker-type and firm-type project factors in
the singular-value representation (T3.21). & Derived. \\
\(\mathbf u_i^\ell,\mathbf v_j^\ell\) & Individual project factors obtained
by lifting rows of \(L_\ell\Sigma_\ell^{1/2}\) and
\(R_\ell\Sigma_\ell^{1/2}\) from types to individuals. & Derived. \\
\(I_{K_\ell}\) & \(K_\ell\)-dimensional identity matrix used as the
interaction coefficient after the SVD factorization. & Standard notation. \\
\(K\) & An independent draw from the firm marginal \(P_J\) when \(J\) and
\(K\) index a pair of firms in \(H_F\). & LOO-v2, p.\ 13. \\
\end{longtable}
\subsection{3.1 The common empirical interface}

Both approaches begin from
\[
\boxed{\mathcal S_{\mathrm{interface}}
=
\left(m,\ P^{\mathrm{obs}}_{IJ}\right),}
\qquad
m_{ij}=\mathbb E[Y\mid I=i,J=j]
\quad\text{on observed cells}.
\tag{T3.1}
\]
The schedule \(m\) is a wage object; \(P^{\mathrm{obs}}_{IJ}\) is an assignment
object.  The project and BS ask different questions of this same interface.

\begin{enumerate}
\item The project decomposes \(m\) relative to a product reference and asks
  how the observed assignment changes wage dispersion.
\item BS20 maps \((m,P^{\mathrm{obs}})\) into observed-match wage types and
  their correlation.
\item BS24 supplies one structural mechanism that generates both \(m\) and
  \(P^{\mathrm{obs}}\).
\end{enumerate}

\textit{Sources: LOO-v2, Sections 2--8; BS20, pp.\ 5--7; BS24/25,
Sections 2--4.}

\subsection{3.2 Project decomposition versus BS functional}

Let
\[
\mu=\sum_{i,j}p_iq_jm_{ij},
\quad
a_i=\sum_jq_jm_{ij}-\mu,
\quad
b_j=\sum_ip_im_{ij}-\mu,
\tag{T3.2}
\]
and
\[
h_{ij}=m_{ij}-\mu-a_i-b_j.
\tag{T3.3}
\]
Then
\[
m_{ij}=\mu+a_i+b_j+h_{ij},
\tag{T3.4}
\]
with product-weighted restrictions
\[
\sum_ip_ia_i=\sum_jq_jb_j=0,
\quad
\sum_jq_jh_{ij}=0,
\quad
\sum_ip_ih_{ij}=0.
\tag{T3.5}
\]
For example,
\[
\sum_ip_ia_i
=
\sum_{i,j}p_iq_jm_{ij}-\mu
=0,
\tag{T3.5a}
\]
and, using \(\sum_jq_jb_j=0\),
\[
\sum_jq_jh_{ij}
=
\sum_jq_jm_{ij}-\mu-a_i-\sum_jq_jb_j
=
(\mu+a_i)-\mu-a_i
=0.
\tag{T3.5b}
\]
The firm-main-effect and interaction-column restrictions follow by the same
calculation with worker and firm indices interchanged:
\[
\sum_jq_jb_j
=
\sum_{i,j}p_iq_jm_{ij}-\mu
=0,
\qquad
\sum_ip_ih_{ij}
=
(\mu+b_j)-\mu-b_j
=0.
\tag{T3.5c}
\]
These equations define the project objects and show where product-supported
missing cells enter.

By contrast, BS types are
\[
\lambda_i^{\mathrm{BS}}
=
\sum_jP^{\mathrm{obs}}(J=j\mid I=i)m_{ij},
\qquad
\mu_j^{\mathrm{BS}}
=
\sum_iP^{\mathrm{obs}}(I=i\mid J=j)m_{ij}.
\tag{T3.6}
\]
Substituting (T3.4) gives
\[
\lambda_i^{\mathrm{BS}}
=
\mu+a_i
+\mathbb E[b_J+h_{iJ}\mid I=i],
\tag{T3.7}
\]
\[
\mu_j^{\mathrm{BS}}
=
\mu+b_j
+\mathbb E[a_I+h_{Ij}\mid J=j].
\tag{T3.8}
\]
Define
\[
r_i^{\mathrm{comp}}
=\mathbb E[b_J+h_{iJ}\mid I=i],
\qquad
s_j^{\mathrm{comp}}
=\mathbb E[a_I+h_{Ij}\mid J=j].
\tag{T3.8a}
\]
Because \(\sum_i p_i a_i=\sum_jq_jb_j=0\), averaging (T3.7)--(T3.8)
over the observed worker and firm marginals gives
\[
\sum_ip_i\lambda_i^{\mathrm{BS}}
=
\sum_{i,j}w_{ij}m_{ij}
=
\sum_jq_j\mu_j^{\mathrm{BS}}
=\bar w.
\tag{T3.8b}
\]
Subtracting this common observed mean from (T3.7)--(T3.8) gives
\[
\lambda_i^{\mathrm{BS}}-\bar w
=a_i+r_i^{\mathrm{comp}}-\sum_{i'}p_{i'}r_{i'}^{\mathrm{comp}},
\tag{T3.8c}
\]
\[
\mu_j^{\mathrm{BS}}-\bar w
=b_j+s_j^{\mathrm{comp}}-\sum_{j'}q_{j'}s_{j'}^{\mathrm{comp}}.
\tag{T3.8d}
\]
Thus centered BS types equal the project main effects exactly when
\(r_i^{\mathrm{comp}}\) is constant across workers and
\(s_j^{\mathrm{comp}}\) is constant across firms.  Otherwise observed
counterpart composition and interactions enter the BS types.

The BS functional is then
\[
\rho_{\mathrm{BS}}
=
\operatorname{Corr}_{P^{\mathrm{obs}}_{IJ}}
(\lambda_I^{\mathrm{BS}},\mu_J^{\mathrm{BS}}).
\tag{T3.9}
\]
This functional is defined for any observed schedule and assignment law with
finite nonzero type variances.  That functional inclusion does not imply that
the BS20 estimator is valid under every model: its repeated-match sampling
assumptions are an additional layer.

\textit{Sources: LOO-v2, pp.\ 9--13 and 21--23; BS20, pp.\ 5--7.
Equations (T3.7)--(T3.8d) are derived substitutions and centering steps.}

\subsection{3.3 How the later BS model generates the interface}

For type pair \((x,y)\), acceptance and the accepted wage are
\[
\bar z_{x,y}
=
\frac{\bar w_x+\bar\pi_y}{f_{x,y}},
\qquad
p_{x,y}=1-G(\bar z_{x,y}),
\tag{T3.10}
\]
\[
W_{x,y}(z)
=
\bar w_x+\gamma
\bigl(zf_{x,y}-\bar w_x-\bar\pi_y\bigr),
\qquad z\geq\bar z_{x,y}.
\tag{T3.11}
\]
With common meeting and separation rates,
\[
\bar\phi_{x,y}
\propto
u_x^{\mathrm{BS}}v_y^{\mathrm{BS}}p_{x,y},
\qquad
\Pi^{\mathrm{obs}}_{x,y}
=
\frac{\bar\phi_{x,y}}{\sum_{x',y'}\bar\phi_{x',y'}}.
\tag{T3.12}
\]
The accepted cell's mean log wage is
\[
\ell_{x,y}
=
\frac{\displaystyle\int_{\bar z_{x,y}}^\infty
\log W_{x,y}(z)g(z)\,dz}
{1-G(\bar z_{x,y})}.
\tag{T3.13}
\]
Under within-type exchangeability,
\[
m^{\log}_{ij}=\ell_{x_i,y_j}.
\tag{T3.14}
\]
Equations (T3.10)--(T3.14), plus a within-type allocation rule, explicitly
construct \((m^{\log},P^{\mathrm{obs}}_{IJ})\).  Applying (T3.6)--(T3.9) then
produces the BS20 sorting statistic implied by the structural equilibrium.

\textit{Sources: BS24/25 revision, pp.\ 5--9 and 17--18; the interface map is
derived.}

\subsection{3.4 Exact level-wage nesting and the log-wage bridge}

Under the Pareto result used by BS24, the accepted \emph{level} wage satisfies
\[
\mathbb E[W_{x,y}(z)\mid z\geq\bar z_{x,y}]
=
\alpha_x^*+\psi_y^*.
\tag{T3.15}
\]
Under the observed type marginals
\(\pi_x=\sum_y\Pi^{\mathrm{obs}}_{x,y}\) and
\(\kappa_y=\sum_x\Pi^{\mathrm{obs}}_{x,y}\), the exact project decomposition is
\[
\begin{aligned}
\mu^L
&=
\sum_x\pi_x\alpha_x^*
+\sum_y\kappa_y\psi_y^*,\\
a_x^L
&=
\alpha_x^*-\sum_{x'}\pi_{x'}\alpha_{x'}^*,\\
b_y^L
&=
\psi_y^*-\sum_{y'}\kappa_{y'}\psi_{y'}^*,\\
h_{x,y}^L&=0.
\end{aligned}
\tag{T3.16}
\]
Hence the Pareto level-wage schedule is an exact rank-zero project submodel.

The project headline outcome is mean log wage, already defined in (T3.13).
In general,
\[
\ell_{x,y}
=
\mathbb E[\log W_{x,y}(z)\mid z\geq\bar z_{x,y}]
\neq
\log(\alpha_x^*+\psi_y^*),
\tag{T3.17}
\]
because the logarithm and conditional expectation do not commute.  Even in
the deterministic special case where equality with
\(\log(\alpha_x^*+\psi_y^*)\) holds, the logarithm of a sum is not generally
additive.  Double-centering the actual mean-log schedule gives
\[
M^\ell_{x,y}
=
\ell_{x,y}-\ell_{x\cdot}-\ell_{\cdot y}+\ell_{\cdot\cdot}.
\tag{T3.18}
\]
Here
\[
\ell_{x\cdot}=\sum_y\kappa_y\ell_{x,y},
\quad
\ell_{\cdot y}=\sum_x\pi_x\ell_{x,y},
\quad
\ell_{\cdot\cdot}=\sum_{x,y}\pi_x\kappa_y\ell_{x,y}.
\tag{T3.18a}
\]
For finite type supports, the exact project interaction rank required is
\[
K_\ell=\operatorname{rank}(M^\ell).
\tag{T3.19}
\]
Let the thin singular-value decomposition be
\[
M^\ell=L_\ell\Sigma_\ell R_\ell'.
\tag{T3.20}
\]
If the maintained project rank is at least \(K_\ell\), set
\[
\mathbf u_x'=e_x'L_\ell,\qquad
\Lambda=\Sigma_\ell,\qquad
\mathbf v_y'=e_y'R_\ell,
\tag{T3.21}
\]
where \(e_x\) and \(e_y\) are coordinate vectors.  Substitution gives
\(\mathbf u_x'\Lambda\mathbf v_y=M^\ell_{x,y}\) cell by cell, so the
representation is exact.  If the maintained rank is smaller than
\(K_\ell\), truncating the singular-value decomposition supplies a low-rank
approximation, not exact nesting.  Therefore the level-wage theorem cannot be
transferred to log wages without checking (T3.19).

\textit{Source for level-wage additivity: BS24/25 Pareto result; equations
(T3.15)--(T3.21) are derived.}

\subsection{3.5 Semistructural comparison and division of labor}

\begin{longtable}{p{0.19\textwidth}p{0.34\textwidth}p{0.38\textwidth}}
\textbf{Dimension} & \textbf{BS} & \textbf{Project} \\
\hline
Primitive interface & Observed matched wages and assignment. & Same. \\
Main reduction & Worker and firm expected-wage types. & Product-centered
additive and low-rank interaction components. \\
Role of models & BS20 uses models as laboratories; BS24 selects one
search/bargaining mechanism. & Maintains a restricted wage surface while
leaving assignment comparatively flexible. \\
Missing cells & Avoided by changing the target to observed-match types. &
Completed only when explicit low-rank/additive restrictions and the required
support/uniqueness conditions identify the missing cells. \\
Estimator & BS20 repeated-match cross-products. & Project leave-out/completion
procedures for its own functionals. \\
Policy content & BS24 can distinguish ATT and ITT under structural
assumptions. & Descriptive unless an assignment/meeting mechanism is added. \\
\end{longtable}

The recommended division of labor is:

\begin{enumerate}
\item use the project decomposition for the headline wage-surface and
  assignment objects;
\item report \(c_{\mathrm{BS}}\) or \(\rho_{\mathrm{BS}}\) as an
  observed-support benchmark;
\item use BS24 to interpret selection and to warn that accepted-wage
  additivity is not production additivity;
\item make production, meeting, ATT, or ITT claims only after adding the
  corresponding structural block.
\end{enumerate}

\textit{Sources: BS20, Sections 3--4; BS24/25, Sections 2--6; LOO-v2,
Sections 2--8.  The division of labor is derived from the preceding maps.}

\section{Question 4: Does BS extend to Kline and the other nested models?}

\subsection{Task 4 notation}

Tasks 1--3 notation remains unchanged.  The chart defines every new symbol
used in Task 4.  Superscripts identify model-specific objects and prevent
collisions with the project's canonical \((\mu,a_i,b_j,h_{ij})\).

\begin{longtable}{p{0.20\textwidth}p{0.52\textwidth}p{0.21\textwidth}}
\textbf{Symbol} & \textbf{Meaning} & \textbf{Source or status} \\
\hline
\(\mathfrak M\) & A generic candidate economic or statistical model. &
Derived notation. \\
\(\mathbf{1}\{\cdot\}\) & Indicator function, equal to one when its argument
is true and zero otherwise. & Standard notation. \\
\(m^{\mathfrak M}_{ij},P^{\mathfrak M}_{IJ}\) & Conditional mean observed
wage schedule and joint observed-match assignment law generated by
\(\mathfrak M\). & Derived generic interface. \\
\(p_i^{\mathfrak M},q_j^{\mathfrak M}\) & Worker and firm marginals of
\(P^{\mathfrak M}_{IJ}\). & Derived. \\
\(\lambda_i^{\mathrm{BS},\mathfrak M},
\mu_j^{\mathrm{BS},\mathfrak M}\) & BS expected-wage worker and firm types
implied by \(\mathfrak M\).  The indexed, superscripted \(\mu\) is distinct
from the project's scalar grand mean \(\mu\). & Derived translation of
BS20, p.\ 5. \\
\(\rho_{\mathrm{BS}}^{\mathfrak M}\) & BS20 population correlation implied
by model \(\mathfrak M\). & Derived model-specific notation. \\
\(\rho_{\mathrm{BS}}^{\mathrm{data}}\) & BS20 population correlation in the
empirical matched-wage distribution, used as a model-validation target. &
Derived label. \\
\(\mathcal B_W,\mathcal B_G,\mathcal B_E,\mathcal B_S\) & Wage-surface,
assignment/selection, directed-edge, and estimator-sampling blocks in the
proposed project architecture. & Derived synthesis. \\
\(\mathsf M\) & Indicator that a potential meeting/access event becomes an
observed match.  It is unrelated to the interaction matrices called \(M\) in
earlier sections. & Derived selection notation. \\
\(P^{\mathrm{meet}}_{IJ},p^{\mathrm{acc}}_{ij}\) & Distribution of potential
meetings and
conditional acceptance probability
\(p^{\mathrm{acc}}_{ij}
=\Pr(\mathsf M=1\mid I=i,J=j,\text{meeting})\). & Derived generic
selection notation. \\
\(d_{ij}^{\mathrm{obs}}\) & Expected duration, survival, or observation
weight of an accepted \((i,j)\) match.  It is common in a one-shot sample or
when all match-separation rates are equal. & Derived stock-sampling
correction. \\
\(\mathcal G_{ij},\operatorname{ITT}_{ij},\operatorname{ATT}_{ij}\) & Worker gain from
a potential meeting, its mean before acceptance is known, and its conditional
mean given acceptance. & Derived generalization of BS24/25,
pp.\ 9--14. \\
\(D_{it},Y_{it}(j)\) & Firm or sector occupied by worker \(i\) at time \(t\),
and the potential log wage at firm/sector \(j\). & Kline, pp.\ 29--34;
GKLP, pp.\ 12--14. \\
\(X_t,\beta\) & Kline time-varying covariates and their coefficient in the
two-period adjusted wage change. & Kline, pp.\ 15, 18--19. \\
\(R_i,R\) & Covariate-adjusted wage change for mover \(i\), and its vector. &
Kline, pp.\ 18--19. \\
\(\mathcal E_F,|\mathcal E_F|\) & Directed firm-mobility edge set and its
number of edges.  This is distinct from the worker--firm observation graph
\(\mathcal E\) in Task 1. & Kline, pp.\ 16--19; subscript is derived. \\
\(E_{\mathrm{edge}}\) & Mover-by-edge indicator matrix, denoted \(E\) by
Kline. & Kline, p.\ 19; subscript is a clarification. \\
\(B_{\mathrm{inc}}\) & Firm-by-edge incidence matrix, denoted \(B\) by
Kline. & Kline, pp.\ 17--19; subscript is a clarification. \\
\(\Delta,\Delta_{jk}\) & Vector of unrestricted Kline edge effects and the
effect on the directed transition \(j\to k\). & Kline, p.\ 19. \\
\(u^{K}_i,u^K\) & Kline edge-regression residual and its vector.  The
superscript distinguishes it from project interaction factors. & Kline,
p.\ 19; superscript is derived. \\
\(c_{\mathrm{cyc}},\mathcal C,C_{\mathrm{cyc}},\dot\psi,\dot\eta\) & A cycle
vector, a directed cycle, matrix of fundamental cycle directions, projected
firm coefficients, and cycle coefficients in Kline's edge decomposition. &
Kline, pp.\ 18--20; subscripts clarify notation. \\
\(\Delta^{LR}_{jk},\Delta^K_{jk}\) & Edge effect implied by the project
low-rank level surface and an arbitrary target Kline edge effect. & Derived. \\
\(\bar{\mathbf u}_{jk}\) & Mean project worker interaction factor among
workers observed moving from \(j\) to \(k\). & Derived. \\
\(d_u,d_v\) & Dimensions of project worker and firm interaction factors. &
Derived project notation. \\
\(\tau_{jk},\varepsilon_i^{\mathrm{edge}}\) & Unrestricted transition wedge
and mover-level residual in the proposed level-plus-edge supermodel. &
Derived extension. \\
\(\alpha_i,\psi_j\) & Primitive additive project worker and firm effects in
the low-rank wage equation used in Task 4.  The canonical product-weighted
components remain \(a_i,b_j\); they need not equal these primitives when
pair-varying covariates remain. & LOO-v2, pp.\ 21--23. \\
\(\psi_j^{\mathrm{ATE}}\) & Population mean potential-wage level of firm
\(j\), relative to a reference firm, under Kline's no-selection-on-treatment-
effects restriction. & Derived notation from Kline, pp.\ 32--34. \\
\(\alpha_i^A,\psi_j^A,\varepsilon_{it}^A\) & AKM worker effect, firm effect,
and residual. & Kline, p.\ 15; superscripts are clarifications. \\
\(m^A_{ij},\lambda_i^{\mathrm{BS},A},
\mu_j^{\mathrm{BS},A}\) & AKM conditional mean wage and the BS worker and
firm wage types it implies. & Derived. \\
\(\alpha_i^T,\psi_j^T,\beta_0^T\) & Tukey worker coordinate, firm coordinate,
and scalar interaction parameter. & Crippa, Section 2.2. \\
\(m^T_{ij}\) & Tukey conditional mean wage. & Derived. \\
\(\bar\psi_i^{\,T},\bar\alpha_j^{\,T}\) & Observed-assignment conditional
means of the Tukey firm and worker coordinates for worker \(i\) and firm
\(j\). & Derived. \\
\(f_m,\alpha_i^S,\psi_j^S\) & Crippa seriation interaction function and its
scalar worker and firm coordinates. & Crippa, Section 2.3.2. \\
\(m^S_{ij}\) & Conditional mean outcome in the Crippa seriation model. &
Crippa, Section 2.3.2. \\
\(\zeta_{k\ell}\) & BLM static conditional mean wage for worker type \(k\)
and firm class \(\ell\). & Derived relabeling of BLM, pp.\ 702--704. \\
\(\Pi^B_{k\ell},\pi_k^B,\kappa_\ell^B\) & BLM observed type-class match
distribution and its worker-type and firm-class marginals. & Derived. \\
\(\lambda_k^{\mathrm{BS},B},\mu_\ell^{\mathrm{BS},B},
\rho_{\mathrm{BS}}^B\) & BLM type-level BS worker type, firm type, and
sorting correlation. & Derived from BS20 and the BLM interface. \\
\(\bar\zeta^B,c_{\mathrm{BS}}^B,
\sigma_{\lambda,B},\sigma_{\mu,B}\) & BLM matched mean wage, BS covariance,
and standard deviations of the two BLM-implied BS wage types. & Derived in
(T4.16a)--(T4.16b). \\
\(M^B,K_B,L_B\) & Double-centered BLM type-class interaction matrix and
numbers of worker types and firm classes. & Derived from BLM,
pp.\ 702--704. \\
\(r_B,U_B,\Sigma_B,V_B,e_k,e_\ell\) & Rank of \(M^B\), matrices in its thin
singular-value decomposition \(M^B=U_B\Sigma_BV_B'\), and coordinate vectors
selecting worker-type and firm-class rows. & Standard linear algebra; derived
application. \\
\(\zeta_{k\cdot},\zeta_{\cdot\ell},\zeta_{\cdot\cdot}\) & Product-weighted
BLM type-row, class-column, and grand means. & Derived. \\
\(\alpha_k^B,a_\ell^B,b_\ell^B,\bar\alpha^B,\bar a^B,\bar b^B\) &
Worker-type value,
firm-class intercept and slope in BLM's simple regression, and the
corresponding weighted worker-type and slope means. & BLM, p.\ 704; weighted
means are derived. \\
\(g(i),h(j)\) & BLM worker-type and firm-class membership maps. & Derived
selector notation. \\
\(Z,W\) & One-hot worker-type and firm-class membership matrices,
\(Z_{ik}=\mathbf{1}\{g(i)=k\}\) and
\(W_{j\ell}=\mathbf{1}\{h(j)=\ell\}\). & Derived selector notation. \\
\(H,H_{ij},u_i^B,v_j^B,u^B,v^B\) & Individual worker--firm interaction matrix,
its scalar rank-one factors, and the vectors collecting those factors in
\(H=ZM^BW'=u^B(v^B)'\). & Derived from (T4.23)--(T4.23e). \\
\(\zeta\) & Matrix collecting the BLM cell means \(\zeta_{k\ell}\) in the
worked example. & Derived notation. \\
\(X_{it},w_{ijt}\) & GKLP observed worker covariates and wage offered to
worker \(i\) by sector \(j\) at time \(t\). & GKLP, pp.\ 6--14. \\
\(y_{ijt},\psi_{ijt}\) & GKLP output and its log unobserved-productivity
component. & GKLP, pp.\ 6--7. \\
\(A_i^G,\eta_i,\varepsilon_{ijt}^G\) & GKLP general ability,
sector-specific ability, and productivity shock.  GKLP write the first
symbol as \(Z_i\); it is relabeled to preserve \(Z\) for BLM membership. &
GKLP, pp.\ 6--7; relabeling is derived. \\
\(\beta_j,b_j^G,c_j^G,\sigma_\varepsilon^2\) & GKLP sector-specific observed-
skill coefficient, return to unobserved skill, intercept, and productivity-
shock variance. & GKLP, pp.\ 6--13; superscripts clarify collisions. \\
\(m_{i,t-1}^G,\sigma_t^2\) & GKLP posterior mean of sector-specific skill
before period \(t\), and the common posterior variance at experience time
\(t\). & GKLP, pp.\ 8--14; superscript clarifies the project \(m_{ij}\). \\
\(\alpha_{it},\psi_{jt},\mathbf u_{it}\) & Time-indexed additive worker
component, additive sector component, and worker interaction factor used in
the exact fixed-\(t\) GKLP embedding. & Derived in (T4.27). \\
\(\mathcal O_{it}\) & A proposed random sector/job offer or access set used
to augment GKLP with a literal meeting/access margin. & Derived extension. \\
\(m_{ijt}^{G,PI},m_{ijt}^{G,L}\) & GKLP perfect-information and learning
conditional mean log-wage surfaces. & Derived labels for GKLP equations
(9)--(10), pp.\ 12--13. \\
\(x,y,f^{SS}(x,y)\) & Shimer--Smith agent types and deterministic type-pair
flow output. & Shimer--Smith, pp.\ 346--349. \\
\(W^{SS}(x),w^{SS}(x)\) & Shimer--Smith unmatched asset value and its flow
value \(w^{SS}(x)=rW^{SS}(x)\). & Shimer--Smith, pp.\ 348--349. \\
\(S^{SS}(x\mid y),\pi^{SS}(x\mid y)\) & Shimer--Smith personal match surplus
and matched flow payoff to type \(x\). & Shimer--Smith, pp.\ 348--349. \\
\(\mathcal M^{SS}(x),a^{SS}_{xy}\) & Shimer--Smith mutually acceptable type
set and deterministic acceptance indicator. & Shimer--Smith,
pp.\ 347--349. \\
\(\Phi^{SS}_{xy}\) & Stationary observed joint distribution of
Shimer--Smith matched types, after normalization. & Derived from their
steady-state match stocks. \\
\(f^{EK}(x,y),w^{EK}(x,y),c_{\mathrm{search}}\) & Eeckhout--Kircher
production, accepted wage, and search cost. & Eeckhout--Kircher,
pp.\ 882--884; search subscript is a clarification. \\
\(\mathcal A_y^{EK},\mathcal B_x^{EK}\) & Eeckhout--Kircher worker and firm
acceptance sets, relabeled to avoid collision with Kline's incidence matrix. &
Eeckhout--Kircher, p.\ 882. \\
\(a_{\mathrm{EK}},h_{\mathrm{EK}},g_{\mathrm{EK}}\) & Coefficient and
worker- and firm-additive functions in the leading production example
\(f^{EK}(x,y)=a_{\mathrm{EK}}xy+h_{\mathrm{EK}}(x)+g_{\mathrm{EK}}(y)\). &
Eeckhout--Kircher, pp.\ 883--885; subscripts clarify notation. \\
\(\bar x,\bar y,H^{EK}(x,y)\) & Product-reference type means and the
double-centered Eeckhout--Kircher wage interaction. & Derived. \\
\end{longtable}
\subsection{4.1 The common four-part test}

For any model \(\mathfrak M\) that generates an observed match law
\(P^{\mathfrak M}_{IJ}\) and conditional mean wage
\[
m^{\mathfrak M}_{ij}
=
\mathbb E_{\mathfrak M}[Y\mid I=i,J=j,\text{observed match}],
\tag{T4.1}
\]
define
\[
\lambda_i^{\mathrm{BS},\mathfrak M}
=
\sum_jP^{\mathfrak M}(J=j\mid I=i)m^{\mathfrak M}_{ij},
\quad
\mu_j^{\mathrm{BS},\mathfrak M}
=
\sum_iP^{\mathfrak M}(I=i\mid J=j)m^{\mathfrak M}_{ij}.
\tag{T4.2}
\]
Then
\[
\rho_{\mathrm{BS}}^{\mathfrak M}
=
\operatorname{Corr}_{P^{\mathfrak M}_{IJ}}
(\lambda_I^{\mathrm{BS},\mathfrak M},
\mu_J^{\mathrm{BS},\mathfrak M})
\tag{T4.3}
\]
exists whenever the relevant second moments are finite and both type
variances are positive.

For every model one must separately ask:

\begin{enumerate}
\item Does it generate (T4.1)--(T4.3)?
\item Does its observation process imply the BS20 repeated-match estimator
  assumptions?
\item Does it contain potential meetings and acceptance, so that BS24's
  ITT/ATT identity is literal?
\item Which wage, assignment, dynamic, and selection layers are nested in the
  project?
\end{enumerate}

If the model has meetings, acceptance probability
\(p^{\mathrm{acc}}_{ij}\), and gain
\(\mathcal G_{ij}=0\) upon rejection, define
\[
\operatorname{ITT}_{ij}
=\mathbb E[\mathcal G_{ij}\mid i,j,\text{meeting}],
\qquad
\operatorname{ATT}_{ij}
=\mathbb E[\mathcal G_{ij}\mid i,j,\text{meeting, accepted}].
\tag{T4.3a}
\]
The law of total expectation then gives
\[
\operatorname{ITT}_{ij}
=
p^{\mathrm{acc}}_{ij}\operatorname{ATT}_{ij},
\tag{T4.4}
\]
and a cross-sectional match stock satisfies
\[
P^{\mathfrak M}_{IJ}(i,j)
\propto
P^{\mathrm{meet}}_{IJ}(i,j)
p^{\mathrm{acc}}_{ij}d_{ij}^{\mathrm{obs}}.
\tag{T4.5}
\]
Models without these objects cannot use (T4.4) literally.

\textit{Sources: BS20, Sections 2 and 4; BS24/25, pp.\ 9--14.}

\subsection{4.2 Decision table}

\begin{longtable}{p{0.15\textwidth}p{0.18\textwidth}p{0.19\textwidth}p{0.20\textwidth}p{0.19\textwidth}}
\textbf{Model} & \textbf{BS population measure} &
\textbf{BS estimator} & \textbf{Literal selection logic} &
\textbf{Project inclusion} \\
\hline
Kline edges & Not from edges alone; yes after adding wage levels and
assignment. & Additional BS sampling assumptions. & Mover ATT/ATE analogue;
literal ITT needs job access. & Add an edge wedge or restrict edge
composition. \\
AKM/KSS & Yes. & KSS inference does not imply BS assumptions. & None. &
Exact rank zero; KSS is inference. \\
Tukey/Crippa & Yes. & Graph identification does not imply BS independence. &
None. & Tukey exact rank one; seriation requires a rank check. \\
BLM & Yes by type-class plug-in. & Use the fitted-model plug-in or separately
impose BS sampling conditions. & Mobility selection, not meeting versus
acceptance. & Exact static type-class mean layer. \\
GKLP & Yes for a declared sector/time target. & Learning changes worker
states and breaks fixed-type repetition. & Roy sector choice, not meeting
acceptance. & Each time slice has rank-one interaction; learning is dynamic. \\
Shimer--Smith & Yes after orienting the two sides. & Needs a separate panel
sampling layer. & Deterministic acceptance by type pair. & Only low-rank
double-centered payoff surfaces. \\
Eeckhout--Kircher & Yes, but not a sign test for productive sorting. & Long
panels do not by themselves imply BS independence. & Two-stage deterministic
acceptance. & Leading example rank one; full search model outside. \\
\end{longtable}

\textit{Sources: model derivations in Sections 4.3--4.6.}

\subsection{4.3 Kline: edge effects versus the project wage surface}

Kline's unrestricted mover model is
\[
R=E_{\mathrm{edge}}\Delta+u^K,
\qquad
\mathbb E[u^K\mid E_{\mathrm{edge}}]=0.
\tag{T4.6}
\]
Each directed transition \(j\to k\) has its own effect \(\Delta_{jk}\).  AKM
imposes
\[
\Delta=B_{\mathrm{inc}}'\psi,
\tag{T4.7}
\]
which implies zero sums around every mobility cycle.  Indeed, if
\(c_{\mathrm{cyc}}\) is an oriented cycle vector, then
\(B_{\mathrm{inc}}c_{\mathrm{cyc}}=0\), and therefore
\[
c_{\mathrm{cyc}}'\Delta
=
c_{\mathrm{cyc}}'B_{\mathrm{inc}}'\psi
=
(B_{\mathrm{inc}}c_{\mathrm{cyc}})'\psi
=0.
\tag{T4.7a}
\]
Conversely, on a connected graph, orthogonality to every cycle puts
\(\Delta\) in the row space of \(B_{\mathrm{inc}}\), so some \(\psi\)
satisfies (T4.7).  Kline's unrestricted model does not impose this
restriction.

The project's static low-rank level surface is
\[
m_{ij}
=
\alpha_i+\psi_j+\mathbf u_i'\Lambda\mathbf v_j.
\tag{T4.8}
\]
For workers moving from \(j\) to \(k\), let
\[
\bar{\mathbf u}_{jk}
=
\mathbb E[\mathbf u_i\mid D_{i1}=j,D_{i2}=k].
\tag{T4.9}
\]
For an individual mover, the worker additive effect cancels:
\[
m_{ik}-m_{ij}
=
\psi_k-\psi_j
+\mathbf u_i'\Lambda(\mathbf v_k-\mathbf v_j).
\tag{T4.9a}
\]
Taking the conditional expectation of (T4.9a) given
\((D_{i1},D_{i2})=(j,k)\), and then using the definition in (T4.9), gives
\[
\Delta^{LR}_{jk}
=
\psi_k-\psi_j
+\bar{\mathbf u}_{jk}'\Lambda(\mathbf v_k-\mathbf v_j).
\tag{T4.10}
\]
This is a structured edge model.  It is not the unrestricted Kline model for
two reasons:

\begin{enumerate}
\item Kline's \(\Delta_{jk}\) is a selected-mover transition object, whereas
  \(m_{ij}\) is a wage level for a worker--firm pair.
\item Conditional on the mover-composition vectors, (T4.10) shares
  \((\psi,\Lambda,\mathbf v)\) across edges.  A generic unrestricted edge
  table need not admit this factorization.
\end{enumerate}

Treating every \(\bar{\mathbf u}_{jk}\) as a free edge parameter would fit
more edges, but it moves the unrestricted flexibility into an unmodeled
selection/composition block; it does not establish wage-layer nesting.

There are two rigorous routes:

\paragraph{Restrict Kline.}
Kline's no-selection-on-treatment-effects condition,
\[
Y_{i2}(k)-Y_{i2}(j)\perp(D_{i1},D_{i2}),
\tag{T4.11}
\]
makes every edge refer to the same population.  Then
\[
\Delta_{jk}
=
\mathbb E[Y_{i2}(k)-Y_{i2}(j)]
=
\psi_k^{\mathrm{ATE}}-\psi_j^{\mathrm{ATE}},
\tag{T4.12}
\]
where
\(\psi_j^{\mathrm{ATE}}
=\mathbb E[Y_{i2}(j)]-\mathbb E[Y_{i2}(1)]\).
Thus the edges are transitive and lie in the rank-zero AKM subclass.  A weaker
route is to assume directly that mover composition and factors satisfy
(T4.10).

\paragraph{Expand the project.}
Retain the level surface but add an unrestricted transition wedge:
\[
R_i
=
\psi_{D_{i2}}-\psi_{D_{i1}}
+\mathbf u_i'\Lambda
(\mathbf v_{D_{i2}}-\mathbf v_{D_{i1}})
+\tau_{D_{i1},D_{i2}}+\varepsilon_i^{\mathrm{edge}}.
\tag{T4.13}
\]
Impose
\(\mathbb E[\varepsilon_i^{\mathrm{edge}}
\mid D_{i1}=j,D_{i2}=k]=0\).
For any target Kline edge table, set
\[
\tau_{jk}
=
\Delta^K_{jk}
-\left[
\psi_k-\psi_j
+\bar{\mathbf u}_{jk}'\Lambda(\mathbf v_k-\mathbf v_j)
\right].
\tag{T4.14}
\]
Taking the conditional expectation of (T4.13) on edge \(j\to k\), substituting
(T4.14), and using the zero conditional mean of
\(\varepsilon_i^{\mathrm{edge}}\) gives
\[
\mathbb E[R_i\mid D_{i1}=j,D_{i2}=k]
=
\Delta^K_{jk}.
\tag{T4.14a}
\]
Thus (T4.13) reproduces \(\Delta^K\) exactly.  The wedge is a transition,
history, or selection object; if it has nonzero cycle sums, it cannot be
reinterpreted as a time-invariant firm wage level.

\textit{Source: Kline, pp.\ 18--20 and 29--34; equations
(T4.9)--(T4.14) are derived.}

\subsection{4.4 BLM: BS plug-in and exact finite-rank nesting}

Let \(\zeta_{k\ell}\) be the BLM conditional mean wage for worker type \(k\)
and firm class \(\ell\), and let \(\Pi^B_{k\ell}\) be the observed type-class
match probability.  Define its marginals
\[
\pi_k^B=\sum_\ell\Pi^B_{k\ell},
\qquad
\kappa_\ell^B=\sum_k\Pi^B_{k\ell}.
\tag{T4.15}
\]
For type and class cells with positive marginals, the exact BLM-implied BS
types are
\[
\lambda_k^{\mathrm{BS},B}
=
\sum_\ell\frac{\Pi^B_{k\ell}}{\pi_k^B}\zeta_{k\ell},
\qquad
\mu_\ell^{\mathrm{BS},B}
=
\sum_k\frac{\Pi^B_{k\ell}}{\kappa_\ell^B}\zeta_{k\ell}.
\tag{T4.16}
\]
Let
\[
\bar\zeta^B=\sum_{k,\ell}\Pi^B_{k\ell}\zeta_{k\ell}.
\tag{T4.16a}
\]
Then the BLM-implied BS moments are
\[
\begin{aligned}
c_{\mathrm{BS}}^B
&=\sum_{k,\ell}\Pi^B_{k\ell}
(\lambda_k^{\mathrm{BS},B}-\bar\zeta^B)
(\mu_\ell^{\mathrm{BS},B}-\bar\zeta^B),\\
(\sigma_{\lambda,B})^2
&=\sum_k\pi_k^B
(\lambda_k^{\mathrm{BS},B}-\bar\zeta^B)^2,\\
(\sigma_{\mu,B})^2
&=\sum_\ell\kappa_\ell^B
(\mu_\ell^{\mathrm{BS},B}-\bar\zeta^B)^2,\\
\rho_{\mathrm{BS}}^B
&=\frac{c_{\mathrm{BS}}^B}
{\sigma_{\lambda,B}\sigma_{\mu,B}}.
\end{aligned}
\tag{T4.16b}
\]
This is a fitted-model plug-in.  Using the raw BS20 estimator instead still
requires the BS sampling assumptions.

For project nesting, product-double-center the BLM mean surface:
\[
\zeta_{k\cdot}=\sum_\ell\kappa_\ell^B\zeta_{k\ell},
\quad
\zeta_{\cdot\ell}=\sum_k\pi_k^B\zeta_{k\ell},
\quad
\zeta_{\cdot\cdot}
=
\sum_{k,\ell}\pi_k^B\kappa_\ell^B\zeta_{k\ell},
\tag{T4.17}
\]
\[
M^B_{k\ell}
=
\zeta_{k\ell}
-\zeta_{k\cdot}
-\zeta_{\cdot\ell}
+\zeta_{\cdot\cdot}.
\tag{T4.18}
\]
By construction,
\[
\sum_\ell\kappa_\ell^BM^B_{k\ell}=0,
\qquad
\sum_k\pi_k^BM^B_{k\ell}=0.
\tag{T4.19}
\]
Let \(r_B=\operatorname{rank}(M^B)\).  A singular-value decomposition gives
\[
M^B=U_B\Sigma_BV_B',
\tag{T4.20}
\]
so choosing
\[
\mathbf u_k'=e_k'U_B,
\qquad
\Lambda=\Sigma_B,
\qquad
\mathbf v_\ell'=e_\ell'V_B
\tag{T4.21}
\]
produces
\(\mathbf u_k'\Lambda\mathbf v_\ell=M^B_{k\ell}\).  Therefore every finite
BLM static mean table is exactly nested when the maintained project rank is
at least \(r_B\).  This is not circular: \(r_B\) is calculated from the
unrestricted fitted table after double-centering; rank one is a testable
special case, not a presumption.

For BLM's simple regression
\[
\zeta_{k\ell}=a_\ell^B+b_\ell^B\alpha_k^B,
\tag{T4.22}
\]
define the product-weighted means
\[
\bar\alpha^B=\sum_k\pi_k^B\alpha_k^B,
\qquad
\bar a^B=\sum_\ell\kappa_\ell^Ba_\ell^B,
\qquad
\bar b^B=\sum_\ell\kappa_\ell^Bb_\ell^B.
\tag{T4.22a}
\]
The row, column, and grand means in (T4.17) are therefore
\[
\zeta_{k\cdot}=\bar a^B+\bar b^B\alpha_k^B,
\qquad
\zeta_{\cdot\ell}=a_\ell^B+b_\ell^B\bar\alpha^B,
\qquad
\zeta_{\cdot\cdot}=\bar a^B+\bar\alpha^B\bar b^B.
\tag{T4.22b}
\]
Substituting these three expressions and (T4.22) into (T4.18) gives
\[
\begin{aligned}
M^B_{k\ell}
&=
a_\ell^B+b_\ell^B\alpha_k^B
-(\bar a^B+\bar b^B\alpha_k^B)
-(a_\ell^B+b_\ell^B\bar\alpha^B)
+(\bar a^B+\bar\alpha^B\bar b^B)\\
&=
(\alpha_k^B-\bar\alpha^B)(b_\ell^B-\bar b^B).
\end{aligned}
\tag{T4.23}
\]
Thus
\[
M^B
=
\begin{pmatrix}
\alpha_1^B-\bar\alpha^B\\
\vdots\\
\alpha_{K_B}^B-\bar\alpha^B
\end{pmatrix}
\begin{pmatrix}
b_1^B-\bar b^B&\cdots&b_{L_B}^B-\bar b^B
\end{pmatrix},
\tag{T4.23a}
\]
an outer product and hence rank at most one.

To connect the type-class matrix to the individual interaction matrix, let
\(g(i)\) be worker \(i\)'s BLM type and \(h(j)\) firm \(j\)'s BLM class.
Define one-hot membership matrices
\[
Z_{ik}=\mathbf{1}\{g(i)=k\},
\qquad
W_{j\ell}=\mathbf{1}\{h(j)=\ell\},
\tag{T4.23b}
\]
and define \(H_{ij}=M^B_{g(i),h(j)}\).  Then
\[
\begin{aligned}
(ZM^BW')_{ij}
&=\sum_{k=1}^{K_B}\sum_{\ell=1}^{L_B}
Z_{ik}M^B_{k\ell}W_{j\ell}\\
&=M^B_{g(i),h(j)}
=H_{ij}.
\end{aligned}
\tag{T4.23c}
\]
The second equality is not a rank assumption: one-hot membership makes every
term zero except \(k=g(i)\) and \(\ell=h(j)\).  Rank one was established
earlier from (T4.22)--(T4.23).  Now define
\[
u_i^B=\alpha_{g(i)}^B-\bar\alpha^B,
\qquad
v_j^B=b_{h(j)}^B-\bar b^B.
\tag{T4.23d}
\]
Equation (T4.23) then gives
\[
H_{ij}=u_i^Bv_j^B,
\qquad
\boxed{H=ZM^BW'=u^B(v^B)'.}
\tag{T4.23e}
\]
Thus \(ZM^BW'\) is the selector-based lift of the type-class interaction,
whereas \(u^B(v^B)'\) is its rank-one factorization.  The logical order is:
derive \(M^B\), establish its rank, define \(H\), prove the selector identity,
and only then factor \(H\).

\textbf{Worked example.}
Take two equally weighted worker types with
\((\alpha_1^B,\alpha_2^B)=(1,-1)\), and two firm classes with weights
\((\kappa_1^B,\kappa_2^B)=(2/3,1/3)\),
\((a_1^B,a_2^B)=(4,5)\), and
\((b_1^B,b_2^B)=(2,-1)\).  Equation (T4.22) gives
\[
\zeta=
\begin{pmatrix}6&4\\2&6\end{pmatrix},
\qquad
\bar\alpha^B=0,\quad \bar b^B=1.
\tag{T4.23f}
\]
Direct double-centering gives
\[
M^B=
\begin{pmatrix}1&-2\\-1&2\end{pmatrix}
=
\begin{pmatrix}1\\-1\end{pmatrix}
\begin{pmatrix}1&-2\end{pmatrix}.
\tag{T4.23g}
\]
Let the two workers have types \(1,2\), and let three firms have classes
\(1,2,1\).  Then
\[
Z=\begin{pmatrix}1&0\\0&1\end{pmatrix},
\qquad
W=\begin{pmatrix}1&0\\0&1\\1&0\end{pmatrix},
\tag{T4.23h}
\]
so both formulas produce
\[
ZM^BW'
=
\begin{pmatrix}1&-2&1\\-1&2&-1\end{pmatrix}
=
\begin{pmatrix}1\\-1\end{pmatrix}
\begin{pmatrix}1&-2&1\end{pmatrix}
=u^B(v^B)'.
\tag{T4.23i}
\]
Every matrix multiplication and the rank-one factorization are therefore
visible in numbers.  The project nests this static mean layer, not BLM's full
wage distributions, transition process, or classification estimator.

\textit{Source: BLM, pp.\ 702--704; BS mapping and
(T4.17)--(T4.23i) are derived.}

\subsection{4.5 GKLP: static comparative advantage versus learning}

Write \(A_i^G\) for the general-ability term that GKLP denote by \(Z_i\);
this preserves \(Z\) for BLM's worker-type membership matrix.
Under perfect information, GKLP's conditional mean log wage is
\[
m_{ijt}^{G,PI}
=
X_{it}'\beta_j
+A_i^G
+b_j^G\eta_i
+c_j^G
+\frac12(b_j^G)^2\sigma_\varepsilon^2.
\tag{T4.24}
\]
For exact nesting, the project covariate design must include the interactions
\(\mathbf{1}\{J=j\}X_{it}\), so each sector may have its own coefficient
\(\beta_j\).  If the project instead imposes a common covariate coefficient,
this block must be expanded before GKLP is nested.  With those interactions
in the covariate component, set
\[
\alpha_i=A_i^G,
\quad
\psi_j=c_j^G+\frac12(b_j^G)^2\sigma_\varepsilon^2,
\quad
\mathbf u_i=(\eta_i),
\quad
\mathbf v_j=(b_j^G),
\quad
\Lambda=(1).
\tag{T4.25}
\]
Substitution of (T4.25) into the project equation reproduces (T4.24)
exactly.  The perfect-information mean surface is therefore rank one.

With learning,
\[
m_{ijt}^{G,L}
=
X_{it}'\beta_j+A_i^G+c_j^G
+b_j^Gm_{i,t-1}^G
+\frac12(b_j^G)^2\sigma_t^2.
\tag{T4.26}
\]
In GKLP, \(\sigma_t^2\) is the common posterior variance at experience time
\(t\), not a worker-specific state.  Therefore, for a fixed \(t\), define
\[
\alpha_{it}=A_i^G,
\qquad
\psi_{jt}=c_j^G+\frac12(b_j^G)^2\sigma_t^2,
\qquad
\mathbf u_{it}=(m_{i,t-1}^G),
\quad
\mathbf v_j=(b_j^G),
\quad
\Lambda=(1).
\tag{T4.27}
\]
Substitution shows that the only worker-by-sector interaction is
\(b_j^Gm_{i,t-1}^G=\mathbf u_{it}'\Lambda\mathbf v_j\).  Each declared
time slice therefore has interaction rank at most one, not two.  The full
panel is nevertheless not a static worker-factor model because
\(m_{i,t-1}^G\) evolves with publicly observed productivity histories, while
\(\psi_{jt}\) changes with \(t\) through \(\sigma_t^2\).  Exact full-model
nesting requires a dynamic state equation, time-varying additive components,
and the sector-choice law.

The BS population measure can be calculated for a fixed sector/time target.
The BS20 estimator is not automatic across learning states because repeated
wages do not measure a fixed worker type.  GKLP choice is Roy selection, not
a random-meeting/acceptance process; a literal BS ITT requires adding an offer
or access set.

\textit{Source: GKLP, pp.\ 6--14; embeddings are derived.}

\subsection{4.6 Remaining models as explicit corollaries}

\begin{enumerate}
\item \textbf{AKM/KSS.}
  \(m^A_{ij}=\alpha_i^A+\psi_j^A\) is the project's rank-zero case.
  Its BS types still include observed counterpart composition:
  \[
  \lambda_i^{\mathrm{BS},A}
  =\alpha_i^A+\mathbb E[\psi_J^A\mid I=i],
  \quad
  \mu_j^{\mathrm{BS},A}
  =\psi_j^A+\mathbb E[\alpha_I^A\mid J=j].
  \tag{T4.28}
  \]
  KSS changes inference, not the wage surface or the BS sampling assumptions.
\item \textbf{Tukey/Crippa.}
  \[
  m^T_{ij}
  =\alpha_i^T+\psi_j^T+\beta_0^T\alpha_i^T\psi_j^T
  \tag{T4.29}
  \]
  has interaction
  \((\beta_0^T\alpha_i^T)\psi_j^T\), an outer product and therefore rank at
  most one after additive components are removed.  General seriation
  identifies an order, not a completed cardinal low-rank surface; exact
  project inclusion requires checking the centered matrix rank.
\item \textbf{Shimer--Smith.}
  The side-\(x\) matched payoff can be rearranged as
  \[
  \pi^{SS}(x\mid y)
  =\tfrac12f^{SS}(x,y)
  +\tfrac12w^{SS}(x)-\tfrac12w^{SS}(y).
  \tag{T4.30}
  \]
  Double-centering removes \(w^{SS}(x)/2\) and
  \(-w^{SS}(y)/2\), leaving one half of the double-centered production
  surface.  Project inclusion therefore requires that surface to
  have rank no greater than the maintained project rank.  Acceptance is
  deterministic by type pair; a genuine within-pair ITT/ATT wedge requires an
  additional match shock.
\item \textbf{Eeckhout--Kircher.}
  In the leading example,
  \(f^{EK}(x,y)=a_{\mathrm{EK}}xy+h_{\mathrm{EK}}(x)+g_{\mathrm{EK}}(y)\).
  Double-centering gives
  \[
  H^{EK}(x,y)
  =
  a_{\mathrm{EK}}
  \left(xy-x\bar y-\bar x y+\bar x\bar y\right)
  =
  a_{\mathrm{EK}}(x-\bar x)(y-\bar y),
  \tag{T4.31}
  \]
  which is rank one.  The BS wage-type correlation is defined, but their type-
  reversal result means it is not a sign test for productive sorting.  The
  full two-stage search model lies outside the static wage layer.
\end{enumerate}

\textit{Sources: Kline, p.\ 15; Crippa, Sections 2--4; Shimer--Smith,
pp.\ 346--349; Eeckhout--Kircher, pp.\ 882--889.}

\subsection{4.7 Project architecture implied by the audit}

Use four distinct blocks:
\[
\boxed{
\begin{array}{ll}
\mathcal B_W:&\text{static wage surface},\\
\mathcal B_G:&\text{assignment, meetings, acceptance, and duration},\\
\mathcal B_E:&\text{directed mover edges and optional }\tau_{jk},\\
\mathcal B_S:&\text{sampling assumptions for the chosen estimator}.
\end{array}}
\tag{T4.32}
\]
The core paper can estimate \(\mathcal B_W\), report the BS population
functional, and use \(\mathcal B_E\) as a diagnostic.  Activate
\(\mathcal B_G\) only for meeting, ATT, ITT,
production, or policy claims.  Activate an unrestricted edge wedge only if
complete Kline edge inclusion is a stated goal.

\textit{Derived from the model audit.}

\section{Question 5: Exogenous mobility, additive separability, and OLS}

\subsection{Task 5 notation}

Tasks 1--4 notation remains in force.  In particular, \(D_{it}\) is worker
\(i\)'s firm at time \(t\), \(Y_{it}(j)\) is the potential log wage at firm
\(j\), and \((\alpha_i,\psi_j,X_{it},\beta)\) are the worker effect, firm
effect, observed covariates, and covariate coefficients.  The following chart
defines every new symbol.

\begin{longtable}{p{0.20\textwidth}p{0.52\textwidth}p{0.21\textwidth}}
\textbf{Symbol} & \textbf{Meaning} & \textbf{Source or status} \\
\hline
\(\varpi_{ij}\) & Permanent statistical worker--firm match effect.  This
relabels Sorkin--Warwar's \(\mu_{ij}\) to avoid collision with the project
grand mean \(\mu\) and the BS wage symbol \(\omega_{ij}\). &
Sorkin--Warwar, p.\ 6; symbol changed. \\
\(\widetilde\varepsilon_{it}(j)\) & Transitory potential-wage remainder at
firm \(j\), after removing the permanent match component. & Sorkin--Warwar,
pp.\ 6--7; potential-outcome index is added. \\
\(\varepsilon^A_{it}\) & Observed residual from an additive AKM equation:
\(\varepsilon^A_{it}=\varpi_{iD_{it}}+
\widetilde\varepsilon_{it}(D_{it})\). & Derived decomposition. \\
\(\mathcal H_{is},\mathcal H_i\) & The period-\(s\) assignment/covariate cell
\((D_{is},X_{is})\), and the entire assignment/covariate history
\(\{(D_{is},X_{is})\}_{s=1}^T\). & Derived shorthand. \\
\(s,t,T\) & Generic periods and the total number of panel periods. & Standard
panel notation. \\
\(\mathbf x\) & Generic realized covariate vector in the
mobility-conditioning cells. & Derived notation. \\
\(\mathsf A_0\) & Pure additive separability:
\(\varpi_{ij}=0\) for every potential worker--firm pair. & Derived label for
Sorkin--Warwar equation (6) and Volpe's second pillar. \\
\(\mathsf E_\varepsilon\) & Mean exogeneity of the transitory potential-wage
remainder with respect to mobility/covariates. & Derived atomic assumption. \\
\(\mathsf E_\varpi\) & Mean exogeneity of permanent match effects with
respect to mobility/covariates. & Derived atomic assumption. \\
\(\mathsf{EM}_{SW},\mathsf{AS}_{SW}\) & Sorkin--Warwar exogenous mobility and
their composite additive-separability assumption. & Derived labels. \\
\(\mathsf{EM}_{V},\mathsf{AS}_{V}\) & Volpe's mobility pillar and pure
additive-separability pillar. & Derived labels. \\
\(\mathsf{AKM}_{V}\) & Intersection of Volpe's two pillars: mobility
exogeneity and pure additive separability. & Derived label. \\
\(\mathsf{SE}_{K}\) & Kline's strict-exogeneity condition for the statistical
AKM equation. & Kline, pp.\ 14--15. \\
\(g_i(j,k)\) & Worker \(i\)'s permanent conditional-mean wage gain from firm
\(j\) to firm \(k\). & Derived. \\
\(G_{i2}(j,k)\) & Kline's period-2 realized potential-wage gain
\(Y_{i2}(k)-Y_{i2}(j)\). & Derived shorthand for Kline, p.\ 32. \\
\(\mathsf{NSG}_{K}\) & Kline's no-selection-on-treatment-effects assumption:
\(G_{i2}(j,k)\perp(D_{i1},D_{i2})\). & Kline, p.\ 32; label derived. \\
\(X^{\mathrm{OLS}},\beta^{\mathrm{OLS}},
\varepsilon^{\mathrm{OLS}}\) & Generic OLS regressor vector, coefficient
vector, and regression residual in
\(Y=(X^{\mathrm{OLS}})'\beta^{\mathrm{OLS}}
+\varepsilon^{\mathrm{OLS}}\). & Standard notation; superscripts prevent
collisions with project and BS symbols. \\
\(D_i,Y_i,Y_i(0),Y_i(1),\tau_i\) & Binary treatment, observed outcome, two
potential outcomes, and individual treatment effect
\(\tau_i=Y_i(1)-Y_i(0)\).  Binary \(D_i\) is
local to the OLS example and is distinct from panel firm assignment
\(D_{it}\). & Standard causal notation. \\
\(d,d_1,d_2\) & Generic firm choices; \(d_t\) is the firm relevant in period
\(t\) in Kline's exclusion restriction. & Kline, pp.\ 29--30. \\
\(\operatorname{ATE}\) & Average treatment effect
\(\mathbb E[\tau_i]\). & Standard notation. \\
\(\theta_0,\theta_1\) & Intercept and binary-treatment slope in the OLS
representation of the average potential-outcome contrast. & Derived. \\
\(u_i\) & Residual in the saturated binary-treatment regression in
(T5.21a)--(T5.21c). & Derived. \\
\(\Delta A_i\) & First difference \(A_{i2}-A_{i1}\) for any panel variable
\(A\); \(i'\) denotes a second worker. & Standard shorthand. \\
\(\xi_i,e_{it},\mathcal N(0,1)\) & Mean-zero match heterogeneity, a
transitory shock used in the counterexamples, and the standard normal
distribution from which that shock is drawn. & Derived; distribution
notation is standard. \\
\end{longtable}
\subsection{5.1 Common decomposition and atomic assumptions}

Write the potential wage at firm \(j\) as
\[
Y_{it}(j)
=
\alpha_i+\psi_j+X_{it}'\beta
+\varpi_{ij}
+\widetilde\varepsilon_{it}(j),
\tag{T5.1}
\]
with observed wage \(Y_{it}=Y_{it}(D_{it})\).  Here \(\varpi_{ij}\) is the
permanent statistical match effect and
\(\widetilde\varepsilon_{it}(j)\) is the transitory remainder.
If an additive AKM regression omits \(\varpi_{ij}\), its observed residual is
\[
\varepsilon_{it}^A
=
\varpi_{iD_{it}}
+\widetilde\varepsilon_{it}(D_{it}).
\tag{T5.1a}
\]
Thus a mean-zero restriction on \(\varepsilon_{it}^A\) restricts the sum in
(T5.1a); by itself it does not prove that either component is pointwise zero.

Separate three assumptions:
\[
\mathsf A_0:
\quad
\varpi_{ij}=0\quad\forall i,j,
\tag{T5.2}
\]
\[
\mathsf E_\varepsilon:
\quad
\mathbb E[
\widetilde\varepsilon_{it}(j)
\mid D_{is}=d,X_{is}=\mathbf x]=0,
\tag{T5.3}
\]
\[
\mathsf E_\varpi:
\quad
\mathbb E[
\varpi_{ij}
\mid D_{is}=d,X_{is}=\mathbf x]=0.
\tag{T5.4}
\]
Conditions (T5.3)--(T5.4) are required for every admissible
\((i,j,d,s,t,\mathbf x)\).  They are pairwise conditional-mean restrictions.  If
\(\mathcal H_i=\{(D_{is},X_{is}):s=1,\ldots,T\}\) is the full observed
assignment/covariate history, the panel strict-exogeneity condition
\[
\mathbb E[
\widetilde\varepsilon_{it}(D_{it})\mid\mathcal H_i]=0
\tag{T5.4a}
\]
conditions jointly on more information and is generally stronger; the two
forms coincide only under additional restrictions on the data-generating
process.
\(\mathsf A_0\) restricts the wage surface.  The other two restrict selection
into mobility cells.

\textit{Sources: Sorkin--Warwar, pp.\ 6--7; Volpe, PDF p.\ 21; notation
\(\varpi_{ij}\) replaces Sorkin--Warwar's \(\mu_{ij}\) to preserve the
project notation.}

\subsection{5.2 Why Volpe and Sorkin are both correct}

Volpe's two pillars are approximately
\[
\mathsf{AS}_{V}=\mathsf A_0,
\qquad
\mathsf{EM}_{V}=\mathsf E_\varepsilon.
\tag{T5.5}
\]
They are non-nested.

\paragraph{Exogenous mobility without additive separability.}
Let there be two firms,
\[
\varpi_{i0}=\xi_i,
\qquad
\varpi_{i1}=-\xi_i,
\qquad
\mathbb E[\xi_i]=0,
\tag{T5.6}
\]
set transitory errors to zero, take covariates to be constant, and assign
firms independently of \(\xi_i\).
Then \(\mathsf E_\varpi\) and \(\mathsf E_\varepsilon\) hold, but
\(\mathsf A_0\) fails.  The permanent gain
\[
g_i(0,1)=\psi_1-\psi_0-2\xi_i
\tag{T5.7}
\]
is heterogeneous.

\paragraph{Additive separability without exogenous mobility.}
Set \(\varpi_{ij}=0\), take covariates to be constant, let
\(e_{it}\sim\mathcal N(0,1)\), and set
\[
D_{it}=\mathbf{1}\{e_{it}>0\},
\qquad
\widetilde\varepsilon_{it}(j)=e_{it}.
\tag{T5.8}
\]
Then \(\mathsf A_0\) holds, but
\[
\mathbb E[e_{it}\mid D_{it}=1]
=
\sqrt{\frac{2}{\pi}}\neq0,
\tag{T5.9}
\]
so mobility exogeneity fails.

Sorkin--Warwar instead define
\[
\mathsf{EM}_{SW}
=
\mathsf E_\varepsilon\cap\mathsf E_\varpi,
\qquad
\mathsf{AS}_{SW}
=
\mathsf E_\varepsilon\cap\mathsf A_0.
\tag{T5.10}
\]
Because \(\mathsf A_0\) mechanically implies
\(\mathsf E_\varpi\),
\[
\boxed{\mathsf{AS}_{SW}\subset\mathsf{EM}_{SW}.}
\tag{T5.11}
\]
The inclusion is strict by (T5.6)--(T5.7).  Thus Volpe applies
\emph{additive separability} to the pure outcome restriction, while
Sorkin--Warwar apply it
to that restriction bundled with transitory-error exogeneity.

Neither definition requires random matching on the additive worker and firm
effects.  Workers may sort on \(\alpha_i\) and on the vector
\(\{\psi_j:j\in\mathcal J\}\).  The restrictions concern the omitted permanent
interaction and transitory components.  Volpe's full-history AKM condition
can be written
\[
\mathbb E[
\varepsilon_{it}^A
\mid
\alpha_i,\{\psi_{D_{is}},X_{is}\}_{s=1}^T]=0.
\tag{T5.11a}
\]
Under \(\mathsf A_0\), equation (T5.1a) reduces this to a mobility restriction
on the transitory error.  Without \(\mathsf A_0\), it is only a
conditional-mean restriction on the composite residual in (T5.1a).

\textit{Sources: Sorkin--Warwar, pp.\ 6--7; Sorkin Lecture 10,
slides 49/79 and following; Volpe, slide 4/35.  Counterexamples and set
relations are derived.}

\subsection{5.3 Kline's three layers}

\paragraph{Statistical AKM.}
Kline writes
\[
Y_{it}
=
\alpha_i+\psi_{D_{it}}+X_{it}'\beta+\varepsilon^A_{it}
\tag{T5.12}
\]
and assumes
\[
\mathbb E[
\varepsilon^A_{it}
\mid D_{is}=j,X_{is}=\mathbf x]=0
\quad\forall(s,t,j,\mathbf x).
\tag{T5.13}
\]
Kline explains that the maintained AKM equation plus (T5.13) embeds both an
additive expected-wage map and mobility not driven by time-varying wage
innovations.

\paragraph{Causal edge effects.}
For two periods, Kline assumes:
\[
Y_{it}(d_1,d_2)=Y_{it}(d_t)
\quad\text{(exclusion),}
\tag{T5.13a}
\]
\[
\mathbb E[
Y_{i2}(j)-Y_{i1}(j)
\mid D_{i1}=j,D_{i2}=k]=0
\quad\text{(parallel trends),}
\tag{T5.13b}
\]
and
\[
\mathbb E[
Y_{i1}(k)-Y_{i1}(j)
\mid D_{i1}=j,D_{i2}=k]
=
\mathbb E[
Y_{i2}(k)-Y_{i2}(j)
\mid D_{i1}=j,D_{i2}=k]
\equiv\Delta_{jk}
\quad\text{(stationarity of the mover gain).}
\tag{T5.13c}
\]
The causal edge result then follows without an additive wage model:
\[
\begin{aligned}
&\mathbb E[Y_{i2}-Y_{i1}\mid D_{i1}=j,D_{i2}=k]\\
&\overset{\text{exclusion}}=
\mathbb E[Y_{i2}(k)-Y_{i1}(j)\mid D_{i1}=j,D_{i2}=k]\\
&=
\mathbb E[Y_{i2}(k)-Y_{i2}(j)\mid D_{i1}=j,D_{i2}=k]\\
&\quad+
\mathbb E[Y_{i2}(j)-Y_{i1}(j)\mid D_{i1}=j,D_{i2}=k]\\
&\overset{\text{parallel trends}}=
\mathbb E[Y_{i2}(k)-Y_{i2}(j)\mid D_{i1}=j,D_{i2}=k]\\
&\overset{\text{stationarity}}=\Delta_{jk}.
\end{aligned}
\tag{T5.14}
\]
Thus \(\Delta_{jk}\) is the average causal effect for workers selected onto
edge \(j\to k\).  This result allows heterogeneous and nonadditive firm
effects.

\paragraph{Common-population ranking.}
Kline additionally assumes
\[
Y_{i2}(k)-Y_{i2}(j)
\perp
(D_{i1},D_{i2}).
\tag{T5.15}
\]
Then
\[
\Delta_{jk}
=
\mathbb E[Y_{i2}(k)-Y_{i2}(j)]
=
\psi_k^{\mathrm{ATE}}-\psi_j^{\mathrm{ATE}},
\tag{T5.16}
\]
where
\(\psi_j^{\mathrm{ATE}}
=\mathbb E[Y_{i2}(j)]-\mathbb E[Y_{i2}(1)]\).
The last equality is direct:
\[
\psi_k^{\mathrm{ATE}}-\psi_j^{\mathrm{ATE}}
=
\mathbb E[Y_{i2}(k)]-\mathbb E[Y_{i2}(j)]
=
\mathbb E[Y_{i2}(k)-Y_{i2}(j)].
\tag{T5.16a}
\]
Hence all edges refer to the same population and are transitive.  Condition
(T5.15) is not constant effects: heterogeneous gains can remain if their
distribution is independent of mobility.

Sorkin--Warwar impose componentwise conditional \emph{mean} restrictions on
wage levels; Kline (T5.15) imposes full independence of pairwise gains.
Neither set of assumptions generally implies the other without additional
restrictions on levels, differences, and conditioning sets.

\textit{Source: Kline, pp.\ 14--16 and 29--34.}

\subsection{5.4 Relation to ordinary OLS}

In
\[
Y=(X^{\mathrm{OLS}})'\beta^{\mathrm{OLS}}
+\varepsilon^{\mathrm{OLS}},
\tag{T5.17}
\]
\[
\mathbb E[\varepsilon^{\mathrm{OLS}}\mid X^{\mathrm{OLS}}]=0
\quad\Longleftrightarrow\quad
\mathbb E[Y\mid X^{\mathrm{OLS}}]
=(X^{\mathrm{OLS}})'\beta^{\mathrm{OLS}}.
\tag{T5.18}
\]
This is a conditional-mean specification.  It is stronger than the
population projection equation
\(\mathbb E[X^{\mathrm{OLS}}\varepsilon^{\mathrm{OLS}}]=0\), and neither condition is
causal without a treatment definition, consistency, and a selection/design
assumption.

For binary treatment,
\[
Y_i=Y_i(0)+D_i\tau_i,
\qquad
\tau_i=Y_i(1)-Y_i(0).
\tag{T5.19}
\]
If
\[
\mathbb E[Y_i(d)\mid D_i]
=
\mathbb E[Y_i(d)]
\quad(d=0,1),
\tag{T5.20}
\]
then
\[
\begin{aligned}
&\mathbb E[Y_i\mid D_i=1]-\mathbb E[Y_i\mid D_i=0]\\
&=
\mathbb E[Y_i(1)\mid D_i=1]
-\mathbb E[Y_i(0)\mid D_i=0]\\
&=
\mathbb E[Y_i(1)]-\mathbb E[Y_i(0)]\\
&=
\mathbb E[\tau_i]
=\operatorname{ATE},
\end{aligned}
\tag{T5.21}
\]
even when \(\tau_i\) is heterogeneous.  To connect this result directly to
OLS, set
\[
\theta_0=\mathbb E[Y_i(0)],
\qquad
\theta_1=\operatorname{ATE},
\qquad
u_i=Y_i-\theta_0-\theta_1D_i,
\tag{T5.21a}
\]
then substitution of (T5.19) gives
\[
u_i=
Y_i(0)-\theta_0
+D_i(\tau_i-\theta_1).
\tag{T5.21b}
\]
Under (T5.20),
\[
\mathbb E[u_i\mid D_i=0]=0,
\qquad
\mathbb E[u_i\mid D_i=1]=0.
\tag{T5.21c}
\]
The first equality uses mean independence of \(Y_i(0)\); the second uses
mean independence of both potential outcomes, which implies
\(\mathbb E[\tau_i\mid D_i=1]=\mathbb E[\tau_i]=\theta_1\).
Thus the saturated binary-treatment regression satisfies ordinary OLS
conditional-mean exogeneity despite heterogeneous individual effects.

Conversely, if
\(\tau_i=\tau\) for every worker but treatment selects on \(Y_i(0)\), then
\[
\begin{aligned}
&\mathbb E[Y_i\mid D_i=1]
-\mathbb E[Y_i\mid D_i=0]\\
&\qquad =
\tau
+\mathbb E[Y_i(0)\mid D_i=1]
-\mathbb E[Y_i(0)\mid D_i=0],
\end{aligned}
\tag{T5.22}
\]
which is generally not \(\tau\).

Thus homogeneous effects do not imply exogeneity, and exogeneity does not
imply homogeneous effects.  The AKM dictionary is exact:

\begin{longtable}{p{0.30\textwidth}p{0.31\textwidth}p{0.30\textwidth}}
\textbf{OLS concept} & \textbf{AKM analogue} & \textbf{Does not imply} \\
\hline
Correct conditional-mean form & Additive worker-plus-firm expected-wage map &
Mobility exogeneity \\
Mean-independent assignment & Mobility does not select potential wage
components/gains & Constant firm effects \\
Constant treatment effect & No permanent worker--firm interaction &
Mobility exogeneity \\
\(\mathbb E[\varepsilon^{\mathrm{OLS}}\mid X^{\mathrm{OLS}}]=0\) &
Strict exogeneity of the maintained AKM residual &
A structural decomposition explaining why the residual is mean zero \\
\end{longtable}

\textit{Equations (T5.17)--(T5.22) are standard potential-outcome and
conditional-mean derivations.}

\subsection{5.5 Implication for the project}

The project permits
\[
Y_{it}(j)
=
\alpha_i+\psi_j
+\mathbf u_i'\Lambda\mathbf v_j
+X_{it}'\beta
+\widetilde\varepsilon_{it}(j).
\tag{T5.23}
\]
It therefore generally rejects pure additive separability but need not reject
mobility exogeneity.  The final specification must state separately:

\begin{enumerate}
\item the wage-surface restriction on
  \(\mathbf u_i'\Lambda\mathbf v_j\);
\item the mobility restriction on transitory shocks, match components,
  potential wage levels, or gains;
\item the time restrictions needed for causal mover effects; and
\item whether the target is a selected-mover edge effect, an unconditional
  average firm effect, a constant individual effect, or a descriptive
  observed-match statistic.
\end{enumerate}

Low rank does not supply exogenous mobility.  Additivity does not make firm
effects causal.  BS20's descriptive population measure requires neither
assumption merely to exist, although its estimator has its own sampling
conditions.

\textit{Derived from Sections 5.1--5.4 and the project wage equation.}

\end{document}
